\documentclass[11pt]{amsart}
\usepackage[margin=1.15in]{geometry}
\usepackage{amssymb,amsmath,amsthm,mathtools}
\usepackage{needspace}
\usepackage[colorlinks=true,linkcolor=blue,urlcolor=blue,citecolor=blue]{hyperref}

\newtheorem{theorem}{Theorem}
\newtheorem{lemma}[theorem]{Lemma}
\newtheorem{proposition}[theorem]{Proposition}
\newtheorem{corollary}[theorem]{Corollary}
\theoremstyle{definition}
\newtheorem{definition}[theorem]{Definition}
\newtheorem*{thmA}{Theorem A}
\theoremstyle{remark}
\newtheorem{remark}[theorem]{Remark}
\numberwithin{equation}{section}

\newcommand{\linf}{\ell^\infty}
\newcommand{\N}{\mathbb{N}}
\newcommand{\Z}{\mathbb{Z}}
\newcommand{\F}{\mathbb{F}}
\newcommand{\cL}{\mathcal{L}}
\newcommand{\cC}{\mathcal{C}}
\newcommand{\cD}{\mathcal{D}}
\newcommand{\Alt}{\operatorname{Alt}}
\newcommand{\Mult}{\operatorname{Mult}}
\newcommand{\End}{\operatorname{End}}
\newcommand{\coeff}{\operatorname{coeff}}
\newcommand{\ord}{\operatorname{ord}}
\newcommand{\sgn}{\operatorname{sgn}}
\newcommand{\rank}{\operatorname{rank}}
\newcommand{\type}{\operatorname{type}}
\newcommand{\spn}{\operatorname{span}}
\newcommand{\id}{\mathrm{id}}
\newcommand{\sdim}{\operatorname{sdim}}
\newcommand{\incl}{\mathrm{incl}}
\newcommand{\Cop}{C^{\mathrm{op}}}
\newcommand{\Cvec}{C^{\mathrm{vec}}}
\newcommand{\Kh}{\widehat K}
\newcommand{\abs}[1]{\lvert #1\rvert}
\newcommand{\norm}[1]{\lVert #1\rVert}
\setlength{\emergencystretch}{1em}


\title[Alternating precomposition is not analytic in positive characteristic]{Precomposition on continuous alternating maps\\ is not analytic in positive characteristic}
\author{Jack McCarthy}
\thanks{The result, its proof and this write-up were produced with Anthropic's Claude (in Claude Code). The proofs
were adversarially audited by independent AI agents (separate Claude instances), but neither the proofs nor this text
has yet been reviewed by a human expert. Unlike the characteristic-$2$ result of~\cite{M}, Theorem~\ref{thm:main} is
\emph{not} machine-checked in Lean. Of the results of this paper, only the lift criterion
(Proposition~\ref{prop:lift}) and the positive direction of Corollary~\ref{cor:class} are machine-checked, together
with a few auxiliary lemmas and the quoted results of~\cite{PR} and Theorem~A (\S\ref{sec:complements}); the
characteristic-$2$ counterexample of~\cite{M} (a different construction, over complete fields that are not spherically
complete) is machine-checked as well. See \S\ref{sec:lean}.}
% the date is printed under the author, as in the characteristic-2 note
\makeatletter
\g@addto@macro\@setauthors{\vskip6pt{\centering\footnotesize September 2026\par}}
\makeatother

\begin{document}
\begin{abstract}
For normed spaces $E,E',F$ over a nontrivially normed field $K$ and a finite set $\iota$, precomposition
$f\mapsto(m\mapsto m\circ(f,\dots,f))$ sends continuous linear maps $E\to E'$ to continuous linear maps between spaces
of continuous alternating $\iota$\nobreakdash-linear maps. This map is $C^n$ for every \mbox{$n\in\N\cup\{\infty\}$} in
every characteristic~\cite{PR}, and it is analytic when $\abs\iota!$ is invertible in $K$. We show that this condition is
also necessary: for every nontrivially normed field $K$ of characteristic $p>0$ and every $k\ge p$ there are
$K$\nobreakdash-Banach spaces $E$ and $F$ such that, with $E'=E$ and $\abs\iota=k$, precomposition is analytic at no point. So
precomposition is analytic for all Banach spaces over $K$ if and only if $\abs\iota!\neq0$ in $K$. The target $F$
admits no equivalent ultrametric norm. The heart of the proof is an algebraic theorem over finite fields: when $k!=0$, the
family of coordinatewise multiplication operators on sequences has no lift in degree $k$ that satisfies the diagonal
identity and has uniformly bounded support.
The algebraic theorem is proved with Ramsey's theorem and an exact cancellation of sign-weighted clusters of indices,
and is combined with a Banach completion of the exterior power over a Laurent series field, whose constant-coefficient map has bounded
support on bounded sets, and with an ascent to an arbitrary base field through Ingleton's Hahn--Banach theorem.
\end{abstract}
\maketitle

\section{Introduction}\label{sec:intro}
Let $K$ be a nontrivially normed field. For normed spaces $E,E',F$ over $K$ and a finite set $\iota$ we write
$\cL(E,E')$ for the space of continuous linear maps with the operator norm, and $\Alt^\iota(E;F)$ for the normed space
of continuous alternating $\iota$\nobreakdash-linear maps $E^\iota\to F$. \emph{Alternating} is meant in the strong sense: the map
vanishes whenever two of its arguments coincide. We study the precomposition map
\begin{equation}\label{eq:Q}
Q=Q^K_{E,E',F}\colon\cL(E,E')\to\cL\bigl(\Alt^\iota(E';F),\Alt^\iota(E;F)\bigr),\qquad Q(f)(m)=m\circ(f,\dots,f).
\end{equation}
In Mathlib it is \texttt{ContinuousAlternatingMap.\allowbreak{}compContinuousLinearMapCLM}.

\subsection*{Motivation}
In Mathlib, manifolds and vector bundles are modeled on normed spaces (of possibly infinite dimension) over an
arbitrary nontrivially normed field, and their regularity is an element $n\in\N\cup\{\infty,\omega\}$, where $C^\omega$
means analytic. Differential forms need the bundle of continuous alternating maps between two $C^n$ vector bundles to
be a $C^n$ vector bundle again. Its transition functions are those of the given bundles with $\Alt^\iota$ applied, so
this comes down to the regularity of $Q$; compare Lang's notion of a functor of class $C^p$~\cite[Ch.~III, \S4]{La}.
The following is known.
\begin{itemize}
\item Over every $K$, with no hypothesis on the characteristic or the dimension, $Q$ is of class $C^n$ for every
$n\in\N\cup\{\infty\}$, in a draft pull request to Mathlib~\cite{PR}.
\item If $\abs\iota!$ is invertible in $K$ then $Q$ is analytic, by polarization. In pull requests to Mathlib, this is
done in characteristic~$0$ by Gou\"ezel~\cite{G1}, who uses it to construct the bundle for every $n$ in
characteristic~$0$~\cite{G2}. A proof in every characteristic is recalled in \S\ref{sec:proofs}.
\item If $K$ is ultrametric and $F$ is ultrametric and spherically complete, then $Q$ is analytic for all normed
$E,E'$. A special case, in the setting of~\cite{M}, is stated without proof in~\cite[Remark~16]{M}; a proof of the
general statement has been formalized in Lean, in an unpublished development (\S\ref{sec:lean}); we quote it and do not
prove it. We call it Theorem~A (\S\ref{sec:complements}). It is not used in the proofs of
Theorem~\ref{thm:main} and Corollary~\ref{cor:class}, only in \S\ref{sec:complements}.
\item Over every complete ultrametric field of characteristic $2$ that is not spherically complete, $Q$ is \emph{not}
analytic at $0$ for $E=E'=\linf$, $\abs\iota=2$ and a suitable ultrametric Banach space $F$~\cite{M}. This is
machine-checked in Lean under the hypothesis (K2) of~\cite{M}, a nested sequence of closed balls with empty
intersection, which is equivalent to non-spherical completeness by an elementary argument
(Remark~\ref{rem:K2}). In characteristic $p\ge3$ the case $\abs\iota\ge p$ was left open in~\cite{M} (see
\cite[Remark~18]{M} for $\abs\iota=p$).
\end{itemize}

\subsection*{Results}
We settle the remaining case for Banach spaces whose norms need not be ultrametric, over every base field.

\begin{theorem}[Main Theorem]\label{thm:main}
Let $K$ be a nontrivially normed field of characteristic $p>0$, and let $k\ge p$. There are Banach spaces $E$ and $F$
over $K$ such that, for every finite set $\iota$ with $\abs\iota=k$, the map
\[
Q\colon\cL(E,E)\to\cL\bigl(\Alt^\iota(E;F),\Alt^\iota(E;F)\bigr),\qquad Q(f)(m)=m\circ(f,\dots,f),
\]
is analytic at no point of $\cL(E,E)$. Equivalently, there is no bounded $k$\nobreakdash-linear map $P$ on $\cL(E,E)$ with
$P(f,\dots,f)=Q(f)$ for all $f$. The Banach space $F$ admits no equivalent ultrametric norm.
\end{theorem}

The spaces are explicit (\S\ref{sec:proofs}). Choose $t\in K$ with $0<\abs t<1$, let $\Kh$ be the completion of $K$
and $K_1\subseteq\Kh$ the closure of $\F_p(t)$, which is isometrically isomorphic to the Laurent series field
$\F_p((X))$ with $\abs X=\abs t$. Put $E_1:=\linf(\N,K_1)$, and let $B$ be the completion of the exterior power
$\Lambda^k_{K_1}E_1$ under the projective norm. Then $E$ and $F$ are the completed projective tensor products
$E_1\widehat\otimes_\pi\Kh$ and $B\widehat\otimes_\pi\Kh$ over $K_1$, regarded as $K$\nobreakdash-Banach spaces.

\begin{corollary}\label{cor:class}
Let $K$ be any nontrivially normed field and $\iota$ any finite set. Then $Q^K_{E,E',F}$ is analytic for all
$K$\nobreakdash-Banach spaces $E,E',F$ if and only if $\abs\iota!\ne0$ in $K$. More precisely:
\begin{enumerate}
\item if $\abs\iota!\neq0$ in $K$, then $Q^K_{E,E',F}$ is analytic, indeed given at every point by a finite power
series, for all normed spaces $E,E',F$, complete or not;
\item if $\abs\iota!=0$ in $K$, then there are Banach spaces $E=E'$ and $F$, with $F$ admitting no equivalent
ultrametric norm, such that $Q^K_{E,E,F}$ is analytic at no point.
\end{enumerate}
\end{corollary}

In Mathlib's terms: for the spaces of Theorem~\ref{thm:main}, \texttt{compContinuousLinearMapCLM} satisfies
\texttt{ContDiff K n} for every $n:\N\cup\{\infty\}$ by~\cite{PR} (on the branch of that draft pull request), and \texttt{ContDiffAt K $\omega$} fails at every
point. So the range of exponents in~\cite{PR} is \emph{sharp} over every field of characteristic $p$ as soon as
$\abs\iota\ge p$. By Corollary~\ref{cor:class}, it cannot be extended to $\omega$ for general Banach spaces over any
field in which $\abs\iota!$ vanishes.

By the third item above (Theorem~A, quoted from that unpublished formalization), no counterexample, over any $K$, can
have a spherically complete ultrametric target (a nonzero ultrametric normed space forces $K$ to be ultrametric). Over a
discretely valued base, such as $\F_q((t))$, the target of a Banach counterexample cannot even carry an equivalent
ultrametric norm; we deduce this from Theorem~A in Proposition~\ref{prop:discrete}, so it rests on the same
formalization. The target of
Theorem~\ref{thm:main} has no equivalent ultrametric norm over any base (Proposition~\ref{prop:nonultra}). For
ultrametric targets the question has a different answer, which depends on the base field and is partly open
(\S\ref{sec:complements}).

\subsection*{Strategy}
\emph{Stage 1 (\S\ref{sec:prelim}).} Analyticity of $Q$ at a single point is equivalent to the existence of a bounded
$k$\nobreakdash-linear lift $P$ with $P(f,\dots,f)=Q(f)$: the lift is the $k$\nobreakdash-th Taylor coefficient, identified by uniqueness of
power series in one variable. No division by $k!$ occurs.

\emph{Stage 2 (\S\ref{sec:laurent}).} Over $K_1=\kappa((X))$ with $\kappa$ finite of characteristic $p$, a bounded lift
$P$ for $E_1=\linf(\N,K_1)$ and the target $B$ produces, by evaluation at multiplication operators and at the universal
alternating map $x\mapsto x_1\wedge\dots\wedge x_k$, followed by the constant-coefficient map, a $\kappa$\nobreakdash-multilinear
map
\[
\Psi\colon E_0^k\times E_0^k\to\Lambda^k_\kappa E_0,\qquad E_0=\kappa^\N,
\]
which is alternating in the last $k$ slots, satisfies the identity
\[
\sum_{\sigma\in S_k}\Psi(u_{\sigma(1)},\dots,u_{\sigma(k)};x)=\sum_{\sigma\in S_k}(u_{\sigma(1)}x_1)\wedge\dots\wedge(u_{\sigma(k)}x_k)
\]
(the multilinear instance of the polarized diagonal identity; $\Psi$ satisfies even the full polarized identity,
Lemma~\ref{lem:Psi}), and whose values all have support dimension at most a fixed $d$. The point is that the constant-coefficient map is
$\kappa$\nobreakdash-linear but not $K_1$\nobreakdash-linear, and has bounded support on bounded subsets of $B$ (Proposition~\ref{prop:supp}).

\emph{Stage 3 (\S\ref{sec:FF}).} Over a finite field with $k!=0$ no such $\Psi$ exists (Theorem~\ref{thm:FF}). Ramsey's
theorem makes the coefficients of $\Psi$ depend only on the order pattern of their indices. Sign-weighted clusters of
four consecutive indices then cancel every term that is not attached to the output, a staircase of clusters and the
support bound force the cluster value $\chi(\tau)$, $\tau\in S_k$, to be invariant under adjacent transpositions, and the
diagonal identity gives $\sum_\tau\chi(\tau)=1$. Hence $k!\,\chi=1$, which is impossible. When $k!$ is invertible the
same computation returns $\chi=1/k!$, the value realized by normalized alternatization.

\emph{Stage 4 (\S\ref{sec:ascent}).} Every $\Kh$ of characteristic $p$ contains a closed, spherically complete copy
$K_1$ of $\F_p((X))$. Ingleton's Hahn--Banach theorem~\cite{In} gives a norm-one $K_1$\nobreakdash-linear projection $\Kh\to K_1$.
It makes the projective base change $V\mapsto V\widehat\otimes_\pi\Kh$ well behaved for arbitrary, not necessarily
ultrametric, $K_1$\nobreakdash-Banach spaces $V$, and it lets a bounded lift over $\Kh$ descend to one over $K_1$. An incomplete $K$
is reduced to $\Kh$ by density.

\emph{Stage 5 (\S\ref{sec:proofs}).} Sums of $N$ unit wedges on disjoint blocks of indices have norm between $N/M_r$ and
$N$ in $B$, hence in $F$, where $r=\abs t$ and $M_r:=\max_{l\ge0}(l+1)r^l$ (so $M_r=1$ for $r\le\frac12$). So $F$ has no
equivalent ultrametric norm.

\subsection*{What is machine-checked}
The lift criterion of Stage~1 and the positive direction of Corollary~\ref{cor:class} are formalized in Lean~4
with Mathlib. The $C^n$ statement of~\cite{PR}, the quoted Theorem~A, a few auxiliary lemmas and the characteristic-$2$
counterexample of~\cite{M} (a different construction) are also formalized (\S\ref{sec:lean}). Theorem~\ref{thm:main}
is proved on paper only; it uses Ramsey's and Ingleton's theorems from the literature and is otherwise self-contained.

\subsection*{Organization}
\S\ref{sec:prelim} collects conventions, the lift criterion and facts about exterior powers.
\S\ref{sec:FF} proves the finite-field theorem. \S\ref{sec:laurent} constructs the target over $\kappa((X))$, $\kappa$
finite, and proves the Main Theorem over that field. \S\ref{sec:ascent} carries it to an arbitrary base field. \S\ref{sec:proofs} assembles the
proofs of Theorem~\ref{thm:main} and Corollary~\ref{cor:class}. \S\ref{sec:complements} discusses ultrametric targets,
related results and open problems, and \S\ref{sec:lean} describes the formalization.

\section{Preliminaries}\label{sec:prelim}

\subsection{Conventions}
We write $\N=\{0,1,2,\dots\}$. Throughout, $K$ denotes the base field and $k$ a degree (\cite{M} writes $k$ for
the field). A \emph{normed field} $K$ carries a multiplicative absolute value $\abs{\cdot}$ satisfying the triangle inequality. It is
\emph{nontrivially normed} if some $c\in K$ has $0<\abs c<1$, and \emph{ultrametric} if $\abs{x+y}\le\max(\abs x,\abs y)$.
A \emph{normed $K$\nobreakdash-space} $X$ satisfies $\norm{\lambda x}=\abs\lambda\norm x$; it is \emph{Banach} if it is complete. We
do \emph{not} require norms of normed spaces to be ultrametric. (In non-archimedean functional analysis ``Banach''
often includes ``ultrametric''; for that reading see \S\ref{sec:complements}.) We write $\cL(X,Y)$ for the bounded
linear maps, $\Mult^k(X;Y)$ for the bounded $k$\nobreakdash-linear maps $X^k\to Y$, and $\Alt^k(X;Y)\subseteq\Mult^k(X;Y)$ for the
closed subspace of alternating ones, all with the operator norm; a subscript $K$ records the field when needed. An
alternating map $m$ is antisymmetric: expanding $0=m(\dots,y+z,\dots,y+z,\dots)$ gives
$m(\dots,y,\dots,z,\dots)=-m(\dots,z,\dots,y,\dots)$, and hence
$m(x_{\sigma(1)},\dots,x_{\sigma(k)})=\sgn(\sigma)\,m(x_1,\dots,x_k)$ for $\sigma\in S_k$. In characteristic $2$ the
converse fails.

For a finite set $\iota$ with $\abs\iota=k$, a bijection $\iota\cong\{1,\dots,k\}$ identifies $\Alt^\iota(E;F)$
isometrically with $\Alt^k(E;F)$ and the map~\eqref{eq:Q} with the corresponding map for $\{1,\dots,k\}$; nothing below
depends on the choice. So we work with
\[
Q=Q^K_{E,E',F}\colon\cL(E,E')\to\cL\bigl(\Alt^k(E';F),\Alt^k(E;F)\bigr),
\]
where $Q(f)(m)(x_1,\dots,x_k)=m(fx_1,\dots,fx_k)$.

\begin{definition}\label{def:lift}
A \emph{bounded $k$\nobreakdash-linear lift} of $Q$ is a bounded $k$\nobreakdash-linear map
$P\colon\cL(E,E')^k\to\cL(\Alt^k(E';F),\Alt^k(E;F))$ with $P(f,\dots,f)=Q(f)$ for all $f\in\cL(E,E')$.
\end{definition}

A map $g\colon X\to Z$ between normed spaces is \emph{analytic at} $x_0$ if there are bounded $n$\nobreakdash-linear maps
$p_n\colon X^n\to Z$ ($n\ge0$) and a radius $R>0$ with $\sup_n\norm{p_n}R^n<\infty$ such that
$g(x_0+h)=\sum_{n\ge0}p_n(h,\dots,h)$, a convergent series, whenever $\norm h<R$. Up to shrinking $R$ this is
Mathlib's \texttt{AnalyticAt}~\cite{docs}, and in Mathlib \texttt{ContDiffAt K $\omega$} at a point implies
\texttt{AnalyticAt} there (\texttt{ContDiffAt.analyticAt}).

\begin{lemma}\label{lem:ultra}
A normed field of characteristic $p>0$ is ultrametric.
\end{lemma}

\begin{proof}
For $n\in\N$ the element $n\cdot1$ lies in the prime field $\F_p$, so $(n\cdot1)^p=n\cdot1$ and
$\abs{n\cdot1}\in\{0,1\}$. Hence $\abs{(x+y)^N}\le\sum_{i=0}^N\abs{\tbinom Ni}\abs x^i\abs y^{N-i}\le(N+1)\max(\abs x,\abs
y)^N$ for every $N$. Taking $N$\nobreakdash-th roots and letting $N\to\infty$ gives $\abs{x+y}\le\max(\abs x,\abs y)$.
\end{proof}

\subsection{The lift criterion}

\begin{lemma}[one-variable uniqueness]\label{lem:unique}
Let $K$ be nontrivially normed and $Z$ a normed $K$\nobreakdash-space, not necessarily complete. Let $z_0,z_1,\dots\in Z$, $\rho>0$
and $M\ge0$ with $\norm{z_n}\rho^n\le M$ for all $n$. If $\sum_nt^nz_n$ converges to $0$ for every $t\in K$ with
$\abs t<\rho$, then $z_n=0$ for all $n$.
\end{lemma}

\begin{proof}
Otherwise let $N$ be least with $z_N\neq0$. For $0<\abs t<\rho/2$ the series $\sum_{n>N}t^{n-N}z_n$ converges, being
$t^{-N}$ times the convergent tail $\sum_{n>N}t^nz_n$, and $t^N\bigl(z_N+\sum_{n>N}t^{n-N}z_n\bigr)=0$. Hence
\[
\norm{z_N}=\Bigl\lVert\sum_{n>N}t^{n-N}z_n\Bigr\rVert\le\sum_{n>N}M\rho^{-n}\abs t^{n-N}
=M\rho^{-N}\frac{\abs t/\rho}{1-\abs t/\rho},
\]
where the norm of a convergent series is bounded by the sum of the norms by continuity of the norm on partial sums.
Since $K$ is nontrivially normed, there are $t\ne0$ with $\abs t$ arbitrarily small, so $\norm{z_N}=0$, a
contradiction.
\end{proof}

\begin{proposition}[lift criterion]\label{prop:lift}
Let $K$ be any nontrivially normed field, $E,E',F$ normed $K$\nobreakdash-spaces and $k\ge1$.
\begin{enumerate}
\item If $Q$ has a bounded $k$\nobreakdash-linear lift $P$, then at every $f_0\in\cL(E,E')$ the map $Q$ is given by a finite power
series of infinite radius; in particular $Q$ is analytic at every point.
\item If $Q$ is analytic at some point $f_0$, with power series $(p_n)_n$ there, then $p_k$ is a bounded $k$\nobreakdash-linear
lift of $Q$.
\item Hence the following are equivalent: $Q$ has a bounded $k$\nobreakdash-linear lift; $Q$ is analytic at every point; $Q$ is
analytic at some point. If they fail, then Mathlib's \texttt{ContDiffAt K $\omega$ Q $f_0$} fails at every $f_0$.
\end{enumerate}
\end{proposition}

\begin{proof}
(1) For $y\in\cL(E,E')$, multilinearity gives
\[
Q(f_0+y)=P(f_0+y,\dots,f_0+y)=\sum_{r=0}^kq_r(y,\dots,y),
\]
where $q_r(y_1,\dots,y_r)$ is the sum, over the $r$\nobreakdash-element subsets $S\subseteq\{1,\dots,k\}$, of $P$ evaluated with
$y_1,\dots,y_r$ in the slots of $S$ in increasing order and $f_0$ in the other slots. Each $q_r$ is
bounded $r$\nobreakdash-linear, with $\norm{q_r}\le\binom kr\norm P\norm{f_0}^{k-r}$, so $(q_r)_{r\le k}$ is a power series of
infinite radius for $Q$ at $f_0$.

(2) Let $\incl\colon\Alt^k(E;F)\hookrightarrow\Mult^k(E;F)$ be the isometric inclusion, and put
\[
Z:=\cL\bigl(\Alt^k(E';F),\Mult^k(E;F)\bigr),\qquad \mathcal I(T):=\incl\circ T .
\]
Then $\mathcal I$ is a linear isometry into $Z$, hence injective.
Fix $h\in\cL(E,E')$, $h\neq0$. For $m\in\Alt^k(E';F)$, $x\in E^k$ and $t\in K$, multilinearity of $m$ gives the
\emph{ambient polynomial identity}
\begin{equation}\label{eq:ambient}
\mathcal I\bigl(Q(f_0+th)\bigr)=\sum_{r=0}^kt^rA_r(h),\qquad A_r(h)(m)(x):=\sum_{\abs S=r}m(u^S_1,\dots,u^S_k),
\end{equation}
where $u^S_i=hx_i$ for $i\in S$ and $u^S_i=f_0x_i$ for $i\notin S$. Here $A_r(h)\in Z$ is bounded, with
$\norm{A_r(h)}\le\binom kr\norm h^r\norm{f_0}^{k-r}$, and $A_k(h)=\mathcal I(Q(h))$. (For $r<k$ we do not need to know whether
$A_r(h)(m)$ is alternating; this is why~\eqref{eq:ambient} is written in $Z$.) Let $R$ and $M':=\sup_n\norm{p_n}R^n$ be
as in the definition of analyticity, and put $\rho:=R/\norm h$. For $\abs t<\rho$ we have $\norm{th}<R$, so
$Q(f_0+th)=\sum_nt^np_n(h,\dots,h)$, and applying $\mathcal I$ and~\eqref{eq:ambient} gives $\sum_nt^nz_n=0$ in $Z$, with
$z_n:=\mathcal I(p_n(h,\dots,h))-A_n(h)$ and $A_n(h):=0$ for $n>k$. Since
$\norm{z_n}\rho^n\le M'+\max_{r\le k}\norm{A_r(h)}\rho^r$, Lemma~\ref{lem:unique} gives $z_k=0$, that is,
$\mathcal I(p_k(h,\dots,h))=\mathcal I(Q(h))$, hence $p_k(h,\dots,h)=Q(h)$. For $h=0$ both sides vanish since $k\ge1$.

(3) follows from (1) and (2), and from \texttt{ContDiffAt.analyticAt} for the last clause.
\end{proof}

\begin{remark}\label{rem:liftlean}
Neither proof divides by $k!$. Item (2) says more than that \emph{some} lift exists: every power series of $Q$ at
every point has its degree-$k$ term equal to a lift.
\end{remark}

\subsection{Exterior powers, supports and determinant arrays}\label{ssec:ext}
Let $L$ be a field, $V$ an $L$\nobreakdash-vector space and $k\ge1$. We write $\Lambda^kV=\Lambda^k_LV$ for the algebraic exterior
power, in which $x_1\wedge\dots\wedge x_k=0$ whenever two $x_i$ coincide. We use the standard facts: for a subspace
$W\subseteq V$ the map $\Lambda^kW\to\Lambda^kV$ is injective, and we regard $\Lambda^kW$ as a subspace of $\Lambda^kV$;
if $u_1,\dots,u_d$ is a basis of a finite-dimensional space $U$, the wedges $u_I:=u_{i_1}\wedge\dots\wedge u_{i_k}$,
$I=\{i_1<\dots<i_k\}$, form a basis of $\Lambda^kU$; and every $\omega\in\Lambda^kV$ lies in $\Lambda^kU$ for some
finite-dimensional $U\subseteq V$.

\begin{definition}\label{def:supp}
The \emph{support dimension} of $\omega\in\Lambda^kV$ is
\[
\sdim(\omega):=\min\{\dim W:\ W\subseteq V\text{ finite-dimensional},\ \omega\in\Lambda^kW\}.
\]
\end{definition}

Clearly $\sdim(0)=0$, $\sdim(c\omega)=\sdim(\omega)$ for $c\neq0$, $\sdim(\omega+\omega')\le\sdim(\omega)+\sdim(\omega')$
(use $W+W'$), and $\sdim(\omega)\le\dim\spn S$ if $\omega$ is a sum of wedges of vectors from a finite set $S$. If
$V\subseteq V'$, then $\sdim$ computed in $V'$ is at most $\sdim$ computed in $V$.

Let $S$ be a set and $V\subseteq L^S$ a space of functions on $S$. The \emph{determinant array} is the linear map
\[
\Omega\colon\Lambda^kV\to L^{S^k},\qquad \Omega_{y_1\wedge\dots\wedge y_k}(c_1,\dots,c_k):=\det\bigl(y_b(c_a)\bigr)_{a,b=1}^k .
\]
It is well defined because the right side is multilinear and alternating in $(y_1,\dots,y_k)$ (these are the
columns). For each $\omega$, the function $\Omega_\omega$ is alternating in $(c_1,\dots,c_k)$ (the rows); in particular
$\Omega_\omega(c_{\sigma(1)},\dots,c_{\sigma(k)})=\sgn(\sigma)\,\Omega_\omega(c_1,\dots,c_k)$.

\begin{lemma}\label{lem:detinj}
The determinant array $\Omega\colon\Lambda^kV\to L^{S^k}$ is injective.
\end{lemma}

\begin{proof}
Let $\Omega_\omega=0$ and choose a finite-dimensional $U\subseteq V$ with $\omega\in\Lambda^kU$. The evaluations
$\mathrm{ev}_s|_U$, $s\in S$, separate the points of $U$ (a function in $U$ vanishing at every $s$ is $0$), so they span
the dual space $U^*$. Choose $s_1,\dots,s_d\in S$ such that the $\mathrm{ev}_{s_i}|_U$ form a basis of $U^*$, and let
$u_1,\dots,u_d$ be the dual basis of $U$, so that $u_j(s_i)=\delta_{ij}$. For increasing $k$\nobreakdash-subsets $I,J$ of
$\{1,\dots,d\}$ we get $\Omega_{u_I}(s_{j_1},\dots,s_{j_k})=\det(u_{i_b}(s_{j_a}))_{a,b}=\delta_{IJ}$. Writing
$\omega=\sum_Ic_Iu_I$ therefore gives $0=\Omega_\omega(s_{j_1},\dots,s_{j_k})=c_J$ for every $J$, so $\omega=0$.
\end{proof}

\begin{lemma}[flattening rank is at most the support]\label{lem:flat}
Let $V\subseteq L^S$, $\omega\in\Lambda^kV$ and $\alpha\in\{1,\dots,k\}$. Let $R\subseteq S$ and
$R'\subseteq S^{\{1,\dots,k\}\setminus\{\alpha\}}$ be finite. The $R\times R'$ matrix $M$ with entries
$M(c_\alpha,c'):=\Omega_\omega(c)$, where $c$ has $c_\alpha$ in slot $\alpha$ and the entries of $c'$ in the other
slots, has rank at most $\sdim(\omega)$.
\end{lemma}

\begin{proof}
Write $\omega=\sum_jy_{j,1}\wedge\dots\wedge y_{j,k}$ with all $y_{j,b}$ in a subspace $W$ of dimension $\sdim(\omega)$.
Expanding each determinant along row $\alpha$,
$\det(y_{j,b}(c_a))_{a,b}=\sum_b\pm y_{j,b}(c_\alpha)\cdot\mu_{j,b}(c')$, where the minor $\mu_{j,b}(c')$ does not
involve $c_\alpha$. So for fixed $c'$ the column $c_\alpha\mapsto M(c_\alpha,c')$ is a linear combination of the
restrictions $y_{j,b}|_R$. All columns lie in the restriction of $W$ to $R$, a space of dimension at most
$\sdim(\omega)$.
\end{proof}

\begin{lemma}[contraction]\label{lem:contr}
Let $k\ge2$ and $\varphi=(\varphi_1,\dots,\varphi_{k-1})\in(V^*)^{k-1}$. There is a linear map $c_\varphi\colon\Lambda^kV\to V$
with
\[
c_\varphi(x_1\wedge\dots\wedge x_k)=\sum_{i=1}^k(-1)^{k+i}\det\bigl(\varphi_a(x_b)\bigr)_{1\le a\le k-1,\ b\ne i}\;x_i ,
\]
and $c_\varphi(\Lambda^kW)\subseteq W$ for every subspace $W$. Consequently
$\sdim(\omega)\ge\dim\spn\{c_\varphi(\omega):\varphi\in(V^*)^{k-1}\}$.
\end{lemma}

\begin{proof}
The right side is multilinear in $(x_b)$. For $\psi\in V^*$, applying $\psi$ to it gives, by Laplace expansion along the
last row, the determinant of the $k\times k$ matrix with entries $\varphi_a(x_b)$ in the rows $a<k$ and $\psi(x_b)$ in the last
row; it vanishes when two columns coincide. Since $V^*$ separates the points of $V$, the right side
vanishes whenever $x_a=x_b$ for some $a\ne b$, so it factors through $\Lambda^kV$. Each term is a multiple of some $x_i$, which gives
$c_\varphi(\Lambda^kW)\subseteq W$. If $\omega\in\Lambda^kW$ with $\dim W=\sdim(\omega)$, all $c_\varphi(\omega)$ lie in $W$.
\end{proof}

\section{The multiplier theorem over a finite field}\label{sec:FF}

\subsection{Lifts and three diagonal identities}\label{ssec:diag}
Fix a field $L$, an integer $k\ge1$, $L$\nobreakdash-vector spaces $A$ and $V$, and a linear map $D\colon A\to\End_L(V)$,
$a\mapsto D_a$. A \emph{lift of $D$ in degree $k$} is a $2k$\nobreakdash-linear map $\Psi\colon A^k\times V^k\to\Lambda^kV$ over $L$ that is
alternating in the $V$\nobreakdash-slots. We write $[k]:=\{1,\dots,k\}$. For $m\ge1$, $f\colon[k]\to[m]$, $b=(b_1,\dots,b_m)\in A^m$
and $\lambda\in L^m$ put
\[
b_f:=(b_{f(1)},\dots,b_{f(k)}),\quad \type f:=\bigl(\abs{f^{-1}(1)},\dots,\abs{f^{-1}(m)}\bigr)\in\N^m,\quad
\lambda_f:=\prod_{j=1}^k\lambda_{f(j)}=\lambda^{\type f}.
\]
Three diagonal identities will be compared. The \emph{pointwise identity} is
\begin{equation*}\tag{Pw}\label{eq:Pw}
\Psi(a,\dots,a;x_1,\dots,x_k)=D_ax_1\wedge\dots\wedge D_ax_k\qquad\text{for all }a\in A,\ x\in V^k.
\end{equation*}
The \emph{polarized identity} is: for all $m\ge1$, $b\in A^m$, $\alpha\in\N^m$ with $\abs\alpha=k$, and $x\in V^k$,
\begin{equation*}\tag{Pol}\label{eq:Pol}
\sum_{\type f=\alpha}\Psi(b_f;x)=\sum_{\type f=\alpha}D_{b_{f(1)}}x_1\wedge\dots\wedge D_{b_{f(k)}}x_k .
\end{equation*}
Its \emph{multilinear instance} is: for all $b_1,\dots,b_k\in A$ and $x\in V^k$,
\begin{equation*}\tag{Pol1}\label{eq:Pol1}
\sum_{\sigma\in S_k}\Psi(b_{\sigma(1)},\dots,b_{\sigma(k)};x)
=\sum_{\sigma\in S_k}D_{b_{\sigma(1)}}x_1\wedge\dots\wedge D_{b_{\sigma(k)}}x_k .
\end{equation*}
Here $\sum_{\type f=\alpha}$ runs over the maps $f\colon[k]\to[m]$ of type $\alpha$. \eqref{eq:Pol1} is the instance $m=k$,
$\alpha=(1,\dots,1)$ of \eqref{eq:Pol}, since the maps of type $(1,\dots,1)$ are the permutations. For $b\in A^m$ and
$x\in V^k$ put $a_\lambda:=\sum_i\lambda_ib_i$ and
\begin{align*}
\Delta_{b,x}(\lambda)&:=\Psi(a_\lambda,\dots,a_\lambda;x)-D_{a_\lambda}x_1\wedge\dots\wedge D_{a_\lambda}x_k
=\sum_{\abs\alpha=k}\lambda^\alpha\,\delta_\alpha(b,x),\\
\delta_\alpha(b,x)&:=\sum_{\type f=\alpha}\Bigl[\Psi(b_f;x)-D_{b_{f(1)}}x_1\wedge\dots\wedge D_{b_{f(k)}}x_k\Bigr].
\end{align*}
The expansion uses multilinearity of $\Psi$ in the $A$\nobreakdash-slots, linearity of $D$ and multilinearity of the wedge. So
$\Delta_{b,x}$ is a homogeneous polynomial of degree $k$ in $\lambda=(\lambda_1,\dots,\lambda_m)$ with coefficients in
$\Lambda^kV$: \eqref{eq:Pol} says that every $\Delta_{b,x}$ is the zero polynomial, and \eqref{eq:Pw} says that every $\Delta_{b,x}$
vanishes at every point of $L^m$ (for the converse take $m=1$ and $\lambda=1$).

\begin{lemma}[identity principle]\label{lem:idprinc}
Let $\mathbb K$ be an infinite field, $Y$ a $\mathbb K$\nobreakdash-vector space and $g(\lambda)=\sum_\alpha\lambda^\alpha y_\alpha$ a
polynomial in $\lambda\in\mathbb K^m$ with finitely many nonzero coefficients $y_\alpha\in Y$. If $g(\lambda)=0$ for all
$\lambda\in\mathbb K^m$, then every $y_\alpha=0$.
\end{lemma}

\begin{proof}
The $y_\alpha$ span a finite-dimensional subspace. Its coordinate functionals with respect to a basis turn $g$ into
finitely many scalar polynomials vanishing on $\mathbb K^m$. These are zero, by induction on $m$, since a nonzero
polynomial in one variable has finitely many roots.
\end{proof}

\begin{lemma}\label{lem:forms}
Let $q$ be a prime power, $L'\supseteq\F_q$ a field, $1\le e\le q$, and $g\in L'[\lambda_1,\dots,\lambda_m]$ homogeneous of
degree $e$ with $g(\lambda)=0$ for every $\lambda\in\F_q^m$. Then $g=0$.
\end{lemma}

\begin{proof}
Call a monomial $\lambda^\beta$ \emph{reduced} if every $\beta_i\le q-1$. The indicator function of a point
$a\in\F_q^m$ is $\prod_i\bigl(1-(\lambda_i-a_i)^{q-1}\bigr)$, which expands into reduced monomials. So the $q^m$ reduced
monomials span the $q^m$\nobreakdash-dimensional $L'$\nobreakdash-space of functions $\F_q^m\to L'$, and they are linearly independent as
functions. A monomial of degree $e\le q$ that is not reduced is $\lambda_i^q$ (and then $e=q\ge2$); as a function on
$\F_q^m$ it agrees with the reduced monomial $\lambda_i$, of degree $1<e$. Replacing each $\lambda_i^q$ in $g$ by
$\lambda_i$ therefore sends the monomials of $g$ injectively to reduced monomials and does not change $g$ as a
function. The resulting reduced polynomial vanishes as a function, hence is $0$, and so $g=0$.
\end{proof}

\begin{proposition}[pointwise versus polarized]\label{prop:diag}
Let $\Psi$ be a lift of $D$ in degree $k$ over a field $L$.
\begin{enumerate}
\item \eqref{eq:Pol}$\Rightarrow$\eqref{eq:Pw}$\Rightarrow$\eqref{eq:Pol1}, over every field.
\item If $L$ is infinite, or if $L=\F_q$ and $k\le q$, then \eqref{eq:Pw}$\Rightarrow$\eqref{eq:Pol}.
\item If $L=\F_q$ and $k\ge q+1$, then \eqref{eq:Pw} does not imply \eqref{eq:Pol} in general.
\item If $k!=0$ in $L$, then \eqref{eq:Pol1} does not imply \eqref{eq:Pw} in general.
\end{enumerate}
\end{proposition}

\begin{proof}
(1) \eqref{eq:Pol}$\Rightarrow$\eqref{eq:Pw} is the case $m=1$, $\lambda=1$. For \eqref{eq:Pw}$\Rightarrow$\eqref{eq:Pol1}, let $b_1,\dots,b_k\in A$,
$x\in V^k$, and put $b_S:=\sum_{i\in S}b_i$ for $S\subseteq[k]$. Then
\begin{align*}
\sum_{S\subseteq[k]}(-1)^{k-\abs S}\Psi(b_S,\dots,b_S;x)&=\sum_S(-1)^{k-\abs S}\sum_{f\colon[k]\to S}\Psi(b_f;x)\\
&=\sum_{f\colon[k]\to[k]}\Psi(b_f;x)\sum_{S\supseteq f([k])}(-1)^{k-\abs S}=\sum_{f\text{ bijective}}\Psi(b_f;x),
\end{align*}
because $\sum_{S\supseteq T}(-1)^{k-\abs S}=(1-1)^{k-\abs T}$ is $1$ if $T=[k]$ and $0$ otherwise. The same computation
applies to the multilinear map $(a_1,\dots,a_k)\mapsto D_{a_1}x_1\wedge\dots\wedge D_{a_k}x_k$, whose diagonal is the right
side of \eqref{eq:Pw}. So \eqref{eq:Pw} at the $2^k$ points $b_S$ gives \eqref{eq:Pol1}. All coefficients are integers.

(2) The coefficients of $\Delta_{b,x}$ span a finite-dimensional subspace of $\Lambda^kV$; in a basis of it each
coordinate of $\Delta_{b,x}$ is a scalar form of degree $k$ vanishing on $L^m$. It is zero by Lemma~\ref{lem:idprinc}
if $L$ is infinite, and by Lemma~\ref{lem:forms} if $L=\F_q$ and $k\le q$.

(3) Let $\dim A\ge3$, with $b_1,b_2,b_3\in A$ and functionals $g_1,g_2,g_3$ on $A$ such that
$g_i(b_j)=\delta_{ij}$, and let $y_1,\dots,y_k\in V$ be linearly independent. Put
\[
\theta(a_1,\dots,a_k):=\bigl[g_1(a_1)\cdots g_1(a_q)\,g_2(a_{q+1})-g_1(a_1)\,g_2(a_2)\cdots g_2(a_{q+1})\bigr]
\cdot g_3(a_{q+2})\cdots g_3(a_k)
\]
(for $k=q+1$ the $g_3$\nobreakdash-factor is empty and $\dim A\ge2$ suffices), and $\Theta(a;x):=\theta(a)\,x_1\wedge\dots\wedge x_k$.
Then $\Theta$ is a lift of $D=0$. It satisfies \eqref{eq:Pw}, because
$\theta(a,\dots,a)=(g_1(a)^qg_2(a)-g_1(a)g_2(a)^q)\,g_3(a)^{k-q-1}=0$ since $c^q=c$ for $c\in\F_q$. It
violates \eqref{eq:Pol} for $\alpha=(q,1,k-q-1)$ and $b=(b_1,b_2,b_3)$: the first product contributes to
$\sum_{\type f=\alpha}\theta(b_f)$ exactly for $f=(1,\dots,1,2,3,\dots,3)$, with value $1$; the second product would need
$f(2)=\dots=f(q+1)=2$, that is, at least $q\ge2$ values equal to $2$, which is impossible for type $\alpha$. So the
coefficient $\delta_\alpha(b,y)$ of $\Theta$ is $y_1\wedge\dots\wedge y_k\neq0$.

(4) Take $\bar a\in A$ and a functional $g_1$ with $g_1(\bar a)=1$, and $y_1,\dots,y_k$ linearly independent. The
lift $\Theta'(a;x):=g_1(a_1)\cdots g_1(a_k)\,x_1\wedge\dots\wedge x_k$ of $D=0$ satisfies \eqref{eq:Pol1}, since
$\sum_\sigma\Theta'(b_\sigma;x)=k!\,g_1(b_1)\cdots g_1(b_k)\,x_1\wedge\dots\wedge x_k=0$ for all
$b_1,\dots,b_k\in A$ and $x\in V^k$, but $\Theta'(\bar a,\dots,\bar a;y)=y_1\wedge\dots\wedge y_k\neq0$.
\end{proof}

So over $\F_q$ the pointwise identity at $\F_q$\nobreakdash-points is strictly weaker than the polarized one as soon as $k>q$. The
proof below uses only \eqref{eq:Pol1}, which by Proposition~\ref{prop:diag}(1) is implied by either. The lift built in
\S\ref{sec:laurent} satisfies the full \eqref{eq:Pol} in every degree, because it comes from polarizing over an \emph{infinite}
field (Lemma~\ref{lem:Psi}).

\subsection{The theorem}
From now on in this section $L$ is a \emph{finite} field with $k!=0$ in $L$, that is, $\operatorname{char}L=p\le k$; in
particular $k\ge2$. Put $V:=L^\N$ and let $V_{\mathrm{fin}}:=L^{(\N)}\subseteq V$ be the subspace of finitely supported
sequences. For $u,x\in V$ we write $ux$ for the coordinatewise product; so $D_ux=ux$ is the multiplier family.

\begin{theorem}[multipliers over a finite field]\label{thm:FF}
Let $L$ be a finite field and $k\ge1$ with $k!=0$ in $L$.
\begin{enumerate}
\item There is no $2k$\nobreakdash-linear map $\Psi\colon V_{\mathrm{fin}}^k\times V_{\mathrm{fin}}^k\to\Lambda^k_LV$ over $L$ such that
\begin{enumerate}
\item $\Psi$ is antisymmetric in the last $k$ slots;
\item $\Psi$ satisfies \eqref{eq:Pol1} for $D_ux=ux$: for all $u_1,\dots,u_k,x_1,\dots,x_k\in V_{\mathrm{fin}}$,
\[
\sum_{\sigma\in S_k}\Psi(u_{\sigma(1)},\dots,u_{\sigma(k)};x_1,\dots,x_k)=\sum_{\sigma\in S_k}(u_{\sigma(1)}x_1)\wedge\dots\wedge(u_{\sigma(k)}x_k);
\]
\item for some $d\in\N$, $\sdim(\Psi(u;x))\le d$ for all $u,x\in V_{\mathrm{fin}}^k$ with entries in $\{0,\pm1\}$.
\end{enumerate}
\item Let $E_0=V$ or $E_0=V_{\mathrm{fin}}$. There is no $2k$\nobreakdash-linear map $\Psi\colon E_0^k\times E_0^k\to\Lambda^k_LE_0$ over $L$
that is antisymmetric in the last $k$ slots, satisfies \eqref{eq:Pol1} for all finitely supported $u,x$, and has
$\sdim(\Psi(u;x))\le d$ for all $u,x\in E_0^k$, for some $d$.
\end{enumerate}
\end{theorem}

\emph{(1) implies (2).} Given $\Psi$ as in (2), restrict it to $V_{\mathrm{fin}}^k\times V_{\mathrm{fin}}^k$ and compose with the
injective map $\Lambda^kE_0\to\Lambda^kV$. Antisymmetry and \eqref{eq:Pol1} are preserved, and the support dimension computed
in $V$ is at most the one computed in $E_0$ (\S\ref{ssec:ext}). So the restriction satisfies (a)--(c).

The rest of this section proves (1). Suppose $\Psi$ is as in (1). Write $e_i\in V_{\mathrm{fin}}$ for the unit vectors,
$e_a:=(e_{a_1},\dots,e_{a_k})$ for $a\in\N^k$, and $\Omega$ for the determinant array of \S\ref{ssec:ext} with $S=\N$.

\subsection{Order-homogeneity from Ramsey's theorem}
We use Ramsey's theorem~\cite{Ra} in its infinite form: \emph{for every $n\ge1$ and every coloring of the
$n$\nobreakdash-element subsets of $\N$ with finitely many colors there is an infinite $H\subseteq\N$ all of whose $n$\nobreakdash-element
subsets have the same color.}

The \emph{pattern} of $z\in\N^N$ is the map sending $i$ to the rank of $z_i$ in the set $\{z_1,\dots,z_N\}$. Thus $z$ and
$z'$ have the same pattern if and only if $z_i<z_j\Leftrightarrow z'_i<z'_j$ and $z_i=z_j\Leftrightarrow z'_i=z'_j$ for
all $i,j$.

\begin{lemma}[pattern homogeneity]\label{lem:H}
For every $N\ge1$ and every function $T\colon\N^N\to L$ there is an infinite $H\subseteq\N$ such that $T(z)=T(z')$
whenever $z,z'\in H^N$ have the same pattern.
\end{lemma}

\begin{proof}
Color an $N$\nobreakdash-element set $s=\{s_1<\dots<s_N\}\subseteq\N$ by the vector
\[
\bigl(T(s_{\beta(1)},\dots,s_{\beta(N)})\bigr)_{\beta\colon[N]\to[N]}\in L^{N^N}.
\]
There are finitely many colors because
$L$ is finite; this is the only use of the finiteness of $L$. Ramsey's theorem gives an infinite $H$ on which the color
is constant, say equal to $(t_\beta)_\beta$. Let $z\in H^N$ have value set $\{y_1<\dots<y_r\}$. Pad it with $N-r$
elements of $H$ above $y_r$ to an $N$\nobreakdash-element set $s\subseteq H$; then $s_j=y_j$ for $j\le r$, and
$z=(s_{\beta(1)},\dots,s_{\beta(N)})$, where $\beta(i)$ is the rank of $z_i$ in $\{y_1,\dots,y_r\}$. The map $\beta$
depends only on the pattern of $z$, and $T(z)=t_\beta$.
\end{proof}

\emph{Coefficients.} For $(a;y;c)\in\N^k\times\N^k\times\N^k$ put
\[
T(a;y;c):=\Omega_{\Psi(e_a;e_y)}(c_1,\dots,c_k)\in L .
\]
Lemma~\ref{lem:H} with $N=3k$ gives an infinite set $H=\{h_0<h_1<h_2<\cdots\}$ on which $T$ depends only on the pattern.
For $u_j,v_j\in V_{\mathrm{fin}}$ supported in $H$, multilinearity of $\Psi$ and linearity of $\Omega$ give the finite
expansion
\begin{equation}\label{eq:expand}
\Omega_{\Psi(u;v)}(c)=\sum_{a,y\in H^k}\ \prod_{j=1}^ku_j(a_j)\prod_{j=1}^kv_j(y_j)\cdot T(a;y;c).
\end{equation}
Below, every output index $c$ lies in $H^k$.

\subsection{Clusters and the cluster value}
A \emph{cluster} is a set $C=\{h_n,h_{n+1},h_{n+2},h_{n+3}\}$ of four
\emph{consecutive} elements of $H$; we write $C(i):=h_{n+i-1}$ for $i=1,\dots,4$. Since a cluster is an interval of
$H$, two disjoint clusters are \emph{comparable}: every element of one is smaller than every element of the other; we
write $C<C'$. Clusters need not be adjacent in $H$. Fix the weights
\begin{align*}
s&=(s_1,s_2,s_3,s_4):=(1,0,-1,0)&&\text{on the operator slots},\\
w&=(w_1,w_2,w_3,w_4):=(1,-1,1,-1)&&\text{on the vector slots}.
\end{align*}
Over $\Z$, hence over every field,
\begin{equation*}\tag{W}\label{eq:W}
\sum_is_i=\sum_iw_i=\sum_is_iw_i=\sum_{i<i'}s_iw_{i'}=\sum_{i>i'}s_iw_{i'}=0,\qquad s_1w_1=1.
\end{equation*}
Indeed $\sum_is_i=1-1$, $\sum_iw_i=1-1+1-1$, $s_1w_1=1$, and
\begin{align*}
\sum_is_iw_i&=s_1w_1+s_3w_3=1-1=0,\\
\sum_{i<i'}s_iw_{i'}&=s_1(w_2+w_3+w_4)+s_3w_4=-1+1=0,\\
\sum_{i>i'}s_iw_{i'}&=s_3(w_1+w_2)=-(1-1)=0 .
\end{align*}
For a cluster $C$ put
\[
s_C:=\sum_{i=1}^4s_ie_{C(i)},\qquad w_C:=\sum_{i=1}^4w_ie_{C(i)} .
\]
These are vectors of $V_{\mathrm{fin}}$ with entries in $\{0,\pm1\}$, and $s_Cw_C=e_{C(1)}-e_{C(3)}$.

\emph{The cluster value.} For pairwise disjoint clusters $C_1,\dots,C_k$, one for each slot, define
\begin{align*}
\chi(C_1,\dots,C_k)&:=\Omega_{\Psi(s_{C_1},\dots,s_{C_k};\,w_{C_1},\dots,w_{C_k})}\bigl(C_1(1),\dots,C_k(1)\bigr)\\
&=\sum_{i,i'\in[4]^k}\ \prod_{j=1}^ks_{i_j}w_{i'_j}\cdot
T\bigl(C_1(i_1),\dots,C_k(i_k);\,C_1(i'_1),\dots,C_k(i'_k);\,C_1(1),\dots,C_k(1)\bigr).
\end{align*}
Let $\tau\in S_k$ be the \emph{order pattern} of $(C_1,\dots,C_k)$: $\tau(j)$ is the position of $C_j$ among
$C_1,\dots,C_k$. The pattern of every $3k$\nobreakdash-tuple in the sum is determined by $(i,i')$ and $\tau$: entries in different
clusters compare as the clusters do, and the entries $C_j(i_j)$, $C_j(i'_j)$, $C_j(1)$ in the same cluster compare as
their indices do, including equalities. By Lemma~\ref{lem:H}, $\chi(C_1,\dots,C_k)=:\chi(\tau)$ depends only on $\tau$.

\begin{lemma}[cancellation of output-free clusters]\label{lem:C}
Let $\cC_1,\dots,\cC_k$ be finite families of clusters that are pairwise disjoint as sets of clusters, such that all
clusters in $\bigcup_j\cC_j$ are pairwise disjoint. Put $u_j:=\sum_{C\in\cC_j}s_C$ and $v_j:=\sum_{C\in\cC_j}w_C$. Let
$C_l\in\cC_l$ for $l=1,\dots,k$ and $c:=(C_1(1),\dots,C_k(1))$. Then $\Omega_{\Psi(u;v)}(c)=\chi(C_1,\dots,C_k)$.
\end{lemma}

\begin{proof}
Expand by~\eqref{eq:expand}. Since the clusters are disjoint, a nonzero term has, for each slot $j$, an operator index
$a_j=\Cop_j(i_j)$ and a vector index $y_j=\Cvec_j(i'_j)$ with $\Cop_j,\Cvec_j\in\cC_j$, and weight
$\prod_js_{i_j}w_{i'_j}$. Group the terms according to the \emph{cluster assignment} $(\Cop_j,\Cvec_j)_{j}$.

Call a cluster \emph{free} in a group if it is occupied (it equals some $\Cop_j$ or some $\Cvec_j$) but is none of the output
clusters $C_1,\dots,C_k$. Suppose a group has a free cluster $\Gamma$. Then $\Gamma\in\cC_j$ for exactly one $j$, because
the families are disjoint, so only $a_j$ and/or $y_j$ lie in $\Gamma$; no output index lies in $\Gamma$, because $\Gamma\ne
C_l$ and the clusters are disjoint. Fix all indices outside $\Gamma$ and sum over the positions inside $\Gamma$. Since
$\Gamma$ is an interval of $H$ containing no other entry of the tuple, every other entry is below all of $\Gamma$ or above
all of $\Gamma$, and the pattern of the tuple depends on the positions inside $\Gamma$ only as follows:
\begin{enumerate}
\item[(i)] if only $a_j\in\Gamma$, not at all, and the partial sum is $(\sum_is_i)\cdot T=0$;
\item[(ii)] if only $y_j\in\Gamma$, not at all, and the partial sum is $(\sum_iw_i)\cdot T=0$;
\item[(iii)] if both lie in $\Gamma$, only through whether $i_j<i'_j$, $i_j=i'_j$ or $i_j>i'_j$; with the corresponding
values $T_<,T_=,T_>$ of $T$ the partial sum is
\[
T_<\sum_{i<i'}s_iw_{i'}+T_=\sum_is_iw_i+T_>\sum_{i>i'}s_iw_{i'}=0 .
\]
\end{enumerate}
(The weights of the other slots are a common factor.) All three vanish by~\eqref{eq:W}, so every group with a free
cluster sums to $0$. In a remaining group every occupied cluster is an output cluster. The only output cluster in
$\cC_j$ is $C_j$, so $\Cop_j=\Cvec_j=C_j$ for all $j$. This single group is, by definition, $\chi(C_1,\dots,C_k)$.
\end{proof}

\subsection{The staircase and the diagonal}

\begin{lemma}[staircase]\label{lem:S}
Let $\tau\in S_k$, let $\alpha\neq\beta$ be slots with $\tau(\alpha)=\ell$ and $\tau(\beta)=\ell+1$, and let $\tau'$ be $\tau$
with these two positions exchanged, so $\tau'(\alpha)=\ell+1$, $\tau'(\beta)=\ell$ and $\tau'(\gamma)=\tau(\gamma)$ otherwise.
Then $\chi(\tau)=\chi(\tau')$.
\end{lemma}

\begin{proof}
Let $m\ge1$. Choose pairwise disjoint clusters, in increasing order in $H$: first a cluster $C_\gamma$ for each slot
$\gamma$ with $\tau(\gamma)<\ell$, ordered as $\tau$ prescribes; then $2m$ clusters $A_1<B_1<A_2<B_2<\dots<A_m<B_m$; then a
cluster $C_\gamma$ for each slot $\gamma$ with $\tau(\gamma)>\ell+1$, ordered as $\tau$ prescribes. Put
$\cC_\alpha:=\{A_1,\dots,A_m\}$, $\cC_\beta:=\{B_1,\dots,B_m\}$ and $\cC_\gamma:=\{C_\gamma\}$ for the other slots, and let
$u,v$ be as in Lemma~\ref{lem:C} and $\omega:=\Psi(u;v)$. The entries of $u$ and $v$ lie in $\{0,\pm1\}$, so
$\sdim(\omega)\le d$ by hypothesis~(c) of Theorem~\ref{thm:FF}(1).

Consider the $m\times m$ matrix $M$ with rows indexed by $c_\alpha=A_n(1)$ and columns by $c_\beta=B_{n'}(1)$
($n,n'=1,\dots,m$), all other output indices being fixed at $c_\gamma=C_\gamma(1)$, and entries $M_{nn'}:=\Omega_\omega(c)$.
By Lemma~\ref{lem:C}, applied with the output clusters $A_n\in\cC_\alpha$, $B_{n'}\in\cC_\beta$ and $C_\gamma\in\cC_\gamma$,
$M_{nn'}$ is $\chi$ of the order pattern of $(A_n,B_{n'},(C_\gamma)_\gamma)$. Now $A_n<B_{n'}$ if and only if $n\le n'$; in
that case the pattern is $\tau$, and otherwise it is $\tau'$. Hence
\[
M=\begin{pmatrix}
\chi(\tau)&\chi(\tau)&\cdots&\chi(\tau)\\
\chi(\tau')&\chi(\tau)&\cdots&\chi(\tau)\\
\vdots&\ddots&\ddots&\vdots\\
\chi(\tau')&\cdots&\chi(\tau')&\chi(\tau)
\end{pmatrix}
=\chi(\tau')\,\mathbf 1\mathbf 1^{\mathsf T}+\delta\,U,
\]
with $\chi(\tau)$ on and above the diagonal and $\chi(\tau')$ below it, where $\delta:=\chi(\tau)-\chi(\tau')$,
$\mathbf 1\in L^m$ is the all-ones column, so that $\mathbf 1\mathbf 1^{\mathsf T}$ is the all-ones matrix, and
$U_{nn'}=[n\le n']$.
$M$ is a submatrix of the flattening of $\Omega_\omega$ along slot $\alpha$, so $\rank M\le\sdim(\omega)\le d$ by
Lemma~\ref{lem:flat}. If $\delta\ne0$, then $\delta U$ is invertible ($U$ is upper unitriangular)
and $\rank\mathbf 1\mathbf 1^{\mathsf T}\le1$, so $\rank M\ge m-1$. Taking $m:=d+2$ gives a contradiction, so $\delta=0$.
\end{proof}

\begin{lemma}[the diagonal]\label{lem:D}
$\sum_{\tau\in S_k}\chi(\tau)=1$.
\end{lemma}

\begin{proof}
Take clusters $C_1<\dots<C_k$, and put $u_j:=s_{C_j}$, $v_j:=w_{C_j}$ and
$c:=(C_1(1),\dots,C_k(1))$. Apply $\omega\mapsto\Omega_\omega(c)$ to both sides of~\eqref{eq:Pol1}, with the operator
vectors $u_1,\dots,u_k$ and the vector arguments $v_1,\dots,v_k$.

\emph{Right side.} For $i\ne l$ the vectors $u_i$ and $v_l$ have disjoint supports, so $u_iv_l=0$ and only $\sigma=\id$
survives. It gives $\Omega_{(u_1v_1)\wedge\dots\wedge(u_kv_k)}(c)=\det\bigl((u_bv_b)(C_a(1))\bigr)_{a,b}$. Since
$u_bv_b=e_{C_b(1)}-e_{C_b(3)}$ and $C_a(1)\ne C_b(3)$ for all $a,b$, this is the determinant of the identity matrix, so
the right side is $1$.

\emph{Left side, one term.} Fix $\sigma\in S_k$ and write $u_\sigma:=(u_{\sigma(1)},\dots,u_{\sigma(k)})$ and
$v_\sigma:=(v_{\sigma(1)},\dots,v_{\sigma(k)})$. Antisymmetry of $\Psi$ in the vector slots gives
$\Psi(u_\sigma;v)=\sgn(\sigma)\Psi(u_\sigma;v_\sigma)$. For every $\omega\in\Lambda^kV$ the function $\Omega_\omega$ is
alternating in its arguments, so $\Omega_\omega(c)=\sgn(\sigma)\,\Omega_\omega(c_{\sigma(1)},\dots,c_{\sigma(k)})$; we
apply this with $\omega=\Psi(u_\sigma;v_\sigma)$. Since $\sgn(\sigma)^2=1$,
\begin{align*}
\Omega_{\Psi(u_\sigma;\,v)}(c)&=\sgn(\sigma)\,\Omega_{\Psi(u_\sigma;\,v_\sigma)}(c)
=\sgn(\sigma)^2\,\Omega_{\Psi(u_\sigma;\,v_\sigma)}(c_{\sigma(1)},\dots,c_{\sigma(k)})\\
&=\Omega_{\Psi(u_\sigma;\,v_\sigma)}(c_{\sigma(1)},\dots,c_{\sigma(k)}).
\end{align*}
In the last expression slot $j$ has operator $s_{C_{\sigma(j)}}$, vector $w_{C_{\sigma(j)}}$ and output
$C_{\sigma(j)}(1)$, so it equals $\chi(C_{\sigma(1)},\dots,C_{\sigma(k)})$ by definition. Its order pattern is $\sigma$,
since $C_{\sigma(j)}$ is the $\sigma(j)$\nobreakdash-th cluster. (Lemma~\ref{lem:C} is not needed here.)

\emph{Left side, summed.} As $\sigma$ runs over $S_k$, so does the order pattern, once each. Comparing with the right
side,
\[
\sum_{\sigma\in S_k}\Omega_{\Psi(u_\sigma;\,v)}(c)=\sum_{\tau\in S_k}\chi(\tau)=1 .\qedhere
\]
\end{proof}

\begin{proof}[Proof of Theorem~\ref{thm:FF}]
The position transpositions $\tau\mapsto(\ell\ \ell{+}1)\circ\tau$, $1\le\ell<k$, generate the action of $S_k$ on itself by
left multiplication, and Lemma~\ref{lem:S} says that $\chi$ is invariant under each of them. So $\chi\equiv\chi_0$ is
constant, and Lemma~\ref{lem:D} gives $k!\,\chi_0=1$ in $L$. Since $k!=0$ in $L$, this reads $0=1$. This proves (1),
and (2) follows as shown above.
\end{proof}

\begin{remark}[where the hypotheses enter]\label{rem:FFhyp}
The proof uses only multilinearity, antisymmetry in the vector slots, the support bound at the finitely many test
inputs of Lemma~\ref{lem:S} (Lemma~\ref{lem:D} does not use the bound), and \eqref{eq:Pol1} at the test inputs of
Lemma~\ref{lem:D}. The finiteness of $L$ enters only in Lemma~\ref{lem:H}, and $k!=0$ only in the last line. By
the inclusion--exclusion in the proof of Proposition~\ref{prop:diag}(1), which uses only multilinearity in the
operator slots, the theorem remains true with the pointwise identity \eqref{eq:Pw} at $L$\nobreakdash-points in place of
\eqref{eq:Pol1}, for this family of multipliers. In characteristic $2$ antisymmetry is the same as symmetry, and the two signs in
Lemma~\ref{lem:D} are trivial.
\end{remark}

\begin{remark}[when $k!$ is invertible]\label{rem:FFunit}
Nothing goes wrong if $k!\ne0$ in $L$. Then the \emph{normalized alternatization}
\[
\Psi_{\mathrm{alt}}(u;x):=\frac1{k!}\sum_{\sigma\in S_k}(u_{\sigma(1)}x_1)\wedge\dots\wedge(u_{\sigma(k)}x_k)
\]
is alternating in $x$, satisfies \eqref{eq:Pol1}, and $\sdim(\Psi_{\mathrm{alt}}(u;x))\le k^2$. Lemmas~\ref{lem:H}--\ref{lem:D}
hold for it, and the computation above returns $\chi\equiv1/k!$: for pairwise disjoint clusters only $\sigma=\id$
survives in $\Psi_{\mathrm{alt}}(s_{C_1},\dots,s_{C_k};w_{C_1},\dots,w_{C_k})$, and the determinant at
$(C_1(1),\dots,C_k(1))$ is that of the identity. The final equation $k!\cdot(1/k!)=1$ is then consistent. So the
obstruction is exactly the vanishing of $k!$, not positive characteristic as such. For example, at $(p,k)=(5,3)$ every
step holds and, since $3!=6=1$ in $\F_5$, the chain stops at $\chi_0=1$.
\end{remark}

\begin{remark}[another possible approach]\label{rem:routeII}
We sketch an alternative way to obtain the input needed in \S\ref{sec:laurent}. Nothing in this paper depends on it.

\emph{The unproved input.} The sketch rests on a statement that we do \emph{not} prove here: the analogue of
Theorem~\ref{thm:FF} over infinite fields $L$ with $k!=0$ in $L$, for lifts that satisfy~\eqref{eq:Pw} and have
uniformly bounded support, of an arbitrary linear family $D\colon A\to\End_L(V)$ whose image modulo the finite-rank
operators is infinite-dimensional.

\emph{Base change and minors.} Let $\Psi$ be a lift over $\F_q$ whose values have support dimension at most $d$.
After the finite base change $\F_q\subseteq\F_{q^N}$, the values of the base change of $\Psi$ at $\F_{q^N}$\nobreakdash-points
have support dimension at most $d_N:=N^{2k}d$, since each input has only $N$ pieces over $\F_q$; choose $N$ with
$q^N>d_N+1$. One would then use that the support dimension of $\omega\in\Lambda^kV$ \emph{equals} the rank of its
contraction matrix, the matrix of coordinates of the vectors $c_\varphi(\omega)$ of Lemma~\ref{lem:contr}, with
$\varphi$ running over tuples of dual basis vectors. Lemma~\ref{lem:contr} gives only the inequality $\ge$, and
\emph{we do not prove the equality}. Granting it, the vanishing of the $(d_N+1)$\nobreakdash-minors of the contraction matrix of a
generic value is a polynomial identity of degree at most $d_N+1<q^N$ in each variable that holds at all
$\F_{q^N}$\nobreakdash-points, so it persists over the algebraic closure $\overline{\F}_q$, and the statement above could be
applied over this infinite field.

\needspace{4\baselineskip}
\emph{What the lift must satisfy.} Three properties of $\Psi$ are needed on the way.
\begin{enumerate}
\item The full identity~\eqref{eq:Pol}, not only~\eqref{eq:Pol1}: for $k>q$ the identity~\eqref{eq:Pw} at
$\F_q$\nobreakdash-points does not survive base change (Proposition~\ref{prop:diag}(3)), whereas~\eqref{eq:Pol}
gives~\eqref{eq:Pw} at every point after base change.
\item Alternation rather than antisymmetry, since alternation is what survives base change in characteristic~$2$.
\item An infinite-dimensional image modulo finite rank. For the multipliers $D_ax=ax$ on $L^\N$ this holds because
the multipliers by the indicators of pairwise disjoint infinite sets are linearly independent modulo finite rank.
\end{enumerate}
The lift of \S\ref{sec:laurent} satisfies~\eqref{eq:Pol} and is alternating (Lemma~\ref{lem:Psi}), and its multiplier
family is the one in~(3), with $L=\kappa$.
\end{remark}

\section{The Banach exterior target over a Laurent series field}\label{sec:laurent}

\subsection{Setting}\label{ssec:laurentsetting}
Throughout this section $\kappa$ is a field, $r\in(0,1)$, and $K_1:=\kappa((X))$ is the field of Laurent series with the
absolute value $\abs a:=r^{\ord a}$ ($\abs0=0$). It is a complete, nontrivially normed, ultrametric field, and $\abs
c=1$ for $c\in\kappa^\times$. For $a\in K_1$ and $n\in\Z$ let $\coeff_n(a)\in\kappa$ be its $X^n$\nobreakdash-coefficient. For a set
$S$ and $y\in\linf(S,K_1)$ we also write
\[
\coeff_0(y):=\bigl(\coeff_0(y(s))\bigr)_{s\in S}\in\kappa^S
\]
for the constant coefficient taken coordinatewise. This \emph{constant-coefficient map} $\coeff_0$ is $\kappa$\nobreakdash-linear but not $K_1$\nobreakdash-linear. Since $\abs a<1$ means $\ord a\ge1$,
\begin{equation}\label{eq:small}
\norm y_\infty<1\ \Longrightarrow\ \coeff_0(y)=0,\qquad\text{hence}\qquad \norm{y-y'}_\infty<1\ \Longrightarrow\ \coeff_0(y)=\coeff_0(y').
\end{equation}
Put $E_1:=\linf(\N,K_1)$ with the sup norm, an ultrametric $K_1$\nobreakdash-Banach space, and let $E_0:=\kappa^\N\subseteq E_1$ be
the $\kappa$\nobreakdash-subspace of sequences with entries in $\kappa$. Every nonzero $z\in E_0$ has $\norm z_\infty=1$. For
$x\in E_1$ and $n\in\Z$ let $x_{[n]}:=(\coeff_n(x(i)))_{i\in\N}\in E_0$ be the $n$\nobreakdash-th Laurent component. For $x\ne0$
the set $\{\ord x(i):x(i)\ne0\}$ is bounded below, because $r^{\ord x(i)}\le\norm x_\infty$, so it has a minimum
$\nu(x)$; then $\norm x_\infty=r^{\nu(x)}$ and $x_{[n]}=0$ for $n<\nu(x)$. Finally put
\[
M_r:=\max_{l\ge0}\,(l+1)r^l .
\]
The maximum exists because $(l+1)r^l\to0$. For $r\le\frac12$ one has $(l+1)r^l\le(l+1)2^{-l}\le1$, so $M_r=1$.

\subsection{The projective exterior norm}\label{ssec:projnorm}
On $\Lambda:=\Lambda^k_{K_1}E_1$ put
\[
\norm\omega_\pi:=\inf\Bigl\{\sum_j\prod_{i=1}^k\norm{x_{j,i}}_\infty\ :\ \omega=\sum_jx_{j,1}\wedge\dots\wedge x_{j,k}\Bigr\},
\]
the infimum over finite decompositions, with an ordinary sum. Let $\Omega^{K_1}\colon\Lambda\to\linf(\N^k,K_1)$ be the determinant
array of \S\ref{ssec:ext} with $L=K_1$, $S=\N$ and $V=E_1$, so that
$\Omega^{K_1}(x_1\wedge\dots\wedge x_k)(i_1,\dots,i_k)=\det(x_b(i_a))_{a,b}$; its values are bounded by the next lemma.

\begin{lemma}\label{lem:proj}
$\norm\cdot_\pi$ is a $K_1$\nobreakdash-norm on $\Lambda$ with $\norm{x_1\wedge\dots\wedge x_k}_\pi\le\prod_i\norm{x_i}_\infty$, and
$\Omega^{K_1}$ is $K_1$\nobreakdash-linear and injective with $\norm{\Omega^{K_1}(\omega)}_\infty\le\norm\omega_\pi$. Let $B$ be the completion of
$(\Lambda,\norm\cdot_\pi)$, a $K_1$\nobreakdash-Banach space, $J\colon B\to\linf(\N^k,K_1)$ the continuous extension of $\Omega^{K_1}$ (so
$\norm J\le1$), and $W_B\colon E_1^k\to B$, $W_B(x):=x_1\wedge\dots\wedge x_k$. Then $W_B\in\Alt^k_{K_1}(E_1;B)$,
$\norm{W_B}\le1$, and $J(W_B(x))=\Omega^{K_1}(x_1\wedge\dots\wedge x_k)$.
\end{lemma}

\begin{proof}
Each entry of $\Omega^{K_1}(x_1\wedge\dots\wedge x_k)$ is a signed sum of products $x_{1}(i_{\sigma(1)})\cdots x_{k}(i_{\sigma(k)})$,
each of absolute value at most $\prod_b\norm{x_b}_\infty$; by the ultrametric inequality in $K_1$ the entry has absolute
value at most $\prod_b\norm{x_b}_\infty$. For a sum of decomposables,
$\norm{\Omega^{K_1}(\omega)}_\infty\le\max_j\prod_i\norm{x_{j,i}}_\infty\le\sum_j\prod_i\norm{x_{j,i}}_\infty$, and taking the infimum
gives $\norm{\Omega^{K_1}(\omega)}_\infty\le\norm\omega_\pi$. Injectivity is Lemma~\ref{lem:detinj} with $L=K_1$, $S=\N$ and
$V=E_1$. For $\lambda\neq0$, rescaling the first factor of every wedge gives $\norm{\lambda\omega}_\pi\le\abs\lambda\norm
\omega_\pi$, and applying this to $\lambda^{-1}$ gives equality; concatenating decompositions gives the triangle
inequality. Positivity: $\norm\omega_\pi\ge\norm{\Omega^{K_1}(\omega)}_\infty>0$ for $\omega\ne0$. Finally $W_B$ is $K_1$\nobreakdash-multilinear,
vanishes when two arguments coincide, and $\norm{W_B(x)}\le\prod_i\norm{x_i}_\infty$.
\end{proof}

$B$ is not ultrametric (Proposition~\ref{prop:nonultra}).

\subsection{The support estimate}\label{ssec:suppest}
Let $\Omega^\kappa\colon\Lambda^k_\kappa E_0\to\kappa^{\N^k}$ be the determinant array over $\kappa$; it is injective by
Lemma~\ref{lem:detinj} with $L=\kappa$ and $V=E_0$. For $y_1,\dots,y_k\in E_0$ we have
$\Omega^{K_1}(y_1\wedge\dots\wedge y_k)=\Omega^\kappa(y_1\wedge\dots\wedge y_k)$: the determinants are the same, with entries in
$\kappa\subseteq K_1$. Support dimensions of elements of $\Lambda^k_\kappa E_0$ are computed in $E_0$ over $\kappa$.

\begin{proposition}[support estimate]\label{prop:supp}
For every $b\in B$ there is a unique $\eta(b)\in\Lambda^k_\kappa E_0$ with $\Omega^\kappa(\eta(b))=\coeff_0(Jb)$. The map $\eta$ is
$\kappa$\nobreakdash-linear, and
\begin{equation}\label{eq:suppest}
\sdim\bigl(\eta(b)\bigr)\le kM_r\,\norm b_B\qquad\text{for all }b\in B .
\end{equation}
\end{proposition}

\begin{proof}
\emph{Step 1: one decomposable.} Let $x_1,\dots,x_k\in E_1$. If some $x_b=0$ the wedge is $0$ and $\eta=0$ works.
Otherwise put $\nu_b:=\nu(x_b)$ and $l:=-\sum_b\nu_b$, so that $\prod_b\norm{x_b}_\infty=r^{-l}$. Fix
$(i_1,\dots,i_k)\in\N^k$ and write $x_b(i_a)=\sum_{n\ge\nu_b}\coeff_n(x_b(i_a))X^n$. Expand
$\det(x_b(i_a))_{a,b}=\sum_\sigma\sgn(\sigma)\prod_bx_b(i_{\sigma(b)})$. The factor $x_b(i_{\sigma(b)})$ has order at
least $\nu_b$, so the constant coefficient of the product is the finite sum
$\sum_n\prod_b\coeff_{n_b}(x_b(i_{\sigma(b)}))$ over
\[
\mathcal N:=\Bigl\{n\in\Z^k:\ n_b\ge\nu_b\text{ for all }b,\ \sum_bn_b=0\Bigr\}.
\]
Exchanging the two finite sums,
\[
\coeff_0\det\bigl(x_b(i_a)\bigr)_{a,b}=\sum_{n\in\mathcal N}\det\bigl(x_{b,[n_b]}(i_a)\bigr)_{a,b},
\]
where $x_{b,[n]}$ is the $n$\nobreakdash-th Laurent component of $x_b$. Hence $\coeff_0\bigl(\Omega^{K_1}(x_1\wedge\dots\wedge x_k)\bigr)=\Omega^\kappa(\eta)$
with
\[
\eta:=\sum_{n\in\mathcal N}x_{1,[n_1]}\wedge\dots\wedge x_{k,[n_k]}\in\Lambda^k_\kappa E_0 .
\]
If $l<0$, then $\sum_bn_b\ge\sum_b\nu_b=-l>0$ for all $n$ with $n_b\ge\nu_b$, so $\mathcal N=\emptyset$ and $\eta=0$: a
decomposable of small norm has zero constant coefficient. If $l\ge0$, every $n\in\mathcal N$ has $n_b-\nu_b\ge0$ and
$\sum_b(n_b-\nu_b)=l$, so $\nu_b\le n_b\le\nu_b+l$. Then $\eta$ is a sum of wedges of vectors from the set
$\{x_{b,[n]}:1\le b\le k,\ \nu_b\le n\le\nu_b+l\}$ of at most $k(l+1)$ vectors, and
\[
\sdim(\eta)\le k(l+1)=k\,(l+1)r^l\cdot r^{-l}\le kM_r\prod_b\norm{x_b}_\infty .
\]
In all cases $\sdim(\eta)\le kM_r\prod_b\norm{x_b}_\infty$.

\emph{Step 2: an element of $\Lambda$.} Let $\omega=\sum_jw_j$ with $w_j=x_{j,1}\wedge\dots\wedge x_{j,k}$, and let
$\eta_j$ be as in Step~1 for $w_j$. Since $\coeff_0$ and $\Omega^{K_1}$ are additive,
$\coeff_0(\Omega^{K_1}(\omega))=\sum_j\coeff_0(\Omega^{K_1}(w_j))=\Omega^\kappa\bigl(\sum_j\eta_j\bigr)$.
By injectivity of $\Omega^\kappa$ the element $\eta(\omega):=\sum_j\eta_j$ does not depend on the decomposition, and
$\sdim(\eta(\omega))\le\sum_j\sdim(\eta_j)\le kM_r\sum_j\prod_i\norm{x_{j,i}}_\infty$. Taking the infimum over
decompositions gives $\sdim(\eta(\omega))\le kM_r\norm\omega_\pi$. (The infimum need not be attained.)

\emph{Step 3: completion.} Let $b\in B$ and $\omega_n\in\Lambda$ with $\norm{b-\omega_n}_B\to0$. Then
$\norm{Jb-\Omega^{K_1}(\omega_n)}_\infty\le\norm{b-\omega_n}_B<1$ for $n\ge n_0$, so by~\eqref{eq:small}
$\coeff_0(Jb)=\coeff_0(\Omega^{K_1}(\omega_n))=\Omega^\kappa(\eta(\omega_n))$ for $n\ge n_0$. Put $\eta(b):=\eta(\omega_n)$, which is independent of
$n\ge n_0$ by injectivity of $\Omega^\kappa$. Then $\sdim(\eta(b))\le kM_r\norm{\omega_n}_\pi$ for all $n\ge n_0$, and
$\norm{\omega_n}_\pi\to\norm b_B$ gives~\eqref{eq:suppest}. Uniqueness is injectivity of $\Omega^\kappa$, and $\eta$ is
$\kappa$\nobreakdash-linear because $J$ is $K_1$\nobreakdash-linear, $\coeff_0$ is $\kappa$\nobreakdash-linear and $\Omega^\kappa$ is injective and linear.
\end{proof}

\subsection{The coefficient lift}\label{ssec:Psi}
Suppose that $P$ is a bounded $k$\nobreakdash-linear lift of $Q^{K_1}_{E_1,E_1,B}$ over $K_1$. For $a\in E_1$ let
$D_a\in\cL(E_1,E_1)$ be coordinatewise multiplication, $D_ax=ax$. The map $a\mapsto D_a$ is $K_1$\nobreakdash-linear and
$\norm{D_a}\le\norm a_\infty$. Define
\begin{align*}
\Phi(a_1,\dots,a_k;x_1,\dots,x_k)&:=P(D_{a_1},\dots,D_{a_k})(W_B)(x_1,\dots,x_k)\in B,\\
\Psi&:=\eta\circ\Phi|_{E_0^{2k}}\colon E_0^k\times E_0^k\to\Lambda^k_\kappa E_0 .
\end{align*}
Since $\Omega^\kappa\circ\eta=\coeff_0\circ J$, the map $\Psi$ is $\coeff_0\circ J\circ\Phi|_{E_0^{2k}}$, read in $\Lambda^k_\kappa E_0$
through the injective map $\Omega^\kappa$.

\begin{lemma}\label{lem:Psi}
The map $\Psi$ has the following properties.
\begin{enumerate}
\item[($\Psi$1)] $\Psi$ is $2k$\nobreakdash-linear over $\kappa$.
\item[($\Psi$2)] $\Psi$ is alternating in the last $k$ slots.
\item[($\Psi$3)] $\sdim(\Psi(z))\le d:=\lfloor kM_r\norm P\rfloor$ for \emph{all} $z\in E_0^{2k}$.
\item[($\Psi$4)] $\Psi$ satisfies \eqref{eq:Pol} for the multiplier family $D^\kappa\colon E_0\to\End_\kappa(E_0)$,
$D^\kappa_ax=ax$: for all $m\ge1$, $b_1,\dots,b_m\in E_0$, $\alpha\in\N^m$ with $\abs\alpha=k$, and $x\in E_0^k$,
\[
\sum_{\type f=\alpha}\Psi(b_f;x)=\sum_{\type f=\alpha}(b_{f(1)}x_1)\wedge\dots\wedge(b_{f(k)}x_k)\qquad\text{in }\Lambda^k_\kappa E_0 .
\]
In particular $\Psi$ satisfies \eqref{eq:Pol1}.
\end{enumerate}
\end{lemma}

\begin{proof}
($\Psi$1) $\Phi$ is $K_1$\nobreakdash-multilinear, hence $\kappa$\nobreakdash-multilinear on the $\kappa$\nobreakdash-subspace $E_0$, and $\eta$ is
$\kappa$\nobreakdash-linear. ($\Psi$2) $P(\dots)(W_B)$ lies in $\Alt^k_{K_1}(E_1;B)$, and $\eta$ is linear. ($\Psi$3) Inputs from
$E_0$ have norm at most $1$, so
$\norm{\Phi(a;x)}\le\norm P\prod_i\norm{D_{a_i}}\,\norm{W_B}\prod_i\norm{x_i}_\infty\le\norm P$; now
use~\eqref{eq:suppest}.

($\Psi$4) Fix $b_1,\dots,b_m\in E_0$ and $x\in E_0^k$, and for $\lambda\in K_1^m$ put $a_\lambda:=\sum_i\lambda_ib_i\in
E_1$. \emph{The diagonal over $K_1$:} the lift property gives, for every $\lambda\in K_1^m$,
\[
\Phi(a_\lambda,\dots,a_\lambda;x)=P(D_{a_\lambda},\dots,D_{a_\lambda})(W_B)(x)=Q(D_{a_\lambda})(W_B)(x)
=W_B(a_\lambda x_1,\dots,a_\lambda x_k).
\]
\emph{Expansion:} expanding both sides by multilinearity gives $\sum_\alpha\lambda^\alpha(L_\alpha-R_\alpha)=0$ for all
$\lambda\in K_1^m$, where
\[
L_\alpha:=\sum_{\type f=\alpha}\Phi(b_f;x)\in B,\qquad R_\alpha:=\sum_{\type f=\alpha}W_B(b_{f(1)}x_1,\dots,b_{f(k)}x_k)\in B .
\]
\emph{Identity principle:} $K_1$ is infinite even when $\kappa$ is finite, so Lemma~\ref{lem:idprinc} in the
$K_1$\nobreakdash-vector space $B$ gives $L_\alpha=R_\alpha$ for every $\alpha$. \emph{Apply $\eta$:} $\eta$ is additive, so
$\eta(L_\alpha)=\sum_{\type f=\alpha}\Psi(b_f;x)$. For $y_1,\dots,y_k\in E_0$, Lemma~\ref{lem:proj} and the remark before
Proposition~\ref{prop:supp} give $J(W_B(y))=\Omega^{K_1}(y_1\wedge\dots\wedge y_k)=\Omega^\kappa(y_1\wedge\dots\wedge y_k)$; this array has
entries in $\kappa$, so $\coeff_0$ fixes it, and uniqueness in Proposition~\ref{prop:supp} gives
$\eta(W_B(y))=y_1\wedge\dots\wedge y_k\in\Lambda^k_\kappa E_0$. With $y_j=b_{f(j)}x_j\in E_0$ this yields
$\eta(R_\alpha)=\sum_{\type f=\alpha}(b_{f(1)}x_1)\wedge\dots\wedge(b_{f(k)}x_k)$.
\end{proof}

The polarization in ($\Psi$4) takes place over the infinite field $K_1$, before the $\kappa$\nobreakdash-linear map $\eta$ is
applied. Evaluating the diagonal only at $\kappa$\nobreakdash-points would give \eqref{eq:Pw} at $\kappa$\nobreakdash-points, which for $\kappa=\F_q$ and
$k>q$ is strictly weaker than \eqref{eq:Pol} (Proposition~\ref{prop:diag}(3)); for the proof below \eqref{eq:Pol1} suffices.

\subsection{No bounded lift over a Laurent series field}

\begin{theorem}\label{thm:laurent}
Let $\kappa$ be a finite field of characteristic $p$, let $k\ge p$ and $r\in(0,1)$, and let $K_1=\kappa((X))$ with
$\abs X=r$. Let $E_1=\linf(\N,K_1)$, and let $B$ be the completion of $\Lambda^k_{K_1}E_1$ under the projective exterior
norm. Then $Q^{K_1}_{E_1,E_1,B}$ admits no bounded $k$\nobreakdash-linear lift over $K_1$. Consequently it is analytic at no point of
$\cL(E_1,E_1)$.
\end{theorem}

\begin{proof}
Suppose $P$ exists, and form $\Psi$ as in \S\ref{ssec:Psi}. By Lemma~\ref{lem:Psi},
$\Psi\colon E_0^k\times E_0^k\to\Lambda^k_\kappa E_0$ with $E_0=\kappa^\N$ is $2k$\nobreakdash-linear over $\kappa$, alternating and hence
antisymmetric in the last $k$ slots, satisfies \eqref{eq:Pol1}, and has $\sdim(\Psi(z))\le d$ for all inputs $z$. Moreover
$k!=0$ in $\kappa$, since $k\ge p=\operatorname{char}\kappa$. This contradicts Theorem~\ref{thm:FF}(2) with $L=\kappa$
and $E_0=\kappa^\N$. The last sentence follows from Proposition~\ref{prop:lift}.
\end{proof}

\begin{remark}\label{rem:infkappa}
Lemmas~\ref{lem:proj} and~\ref{lem:Psi} and Proposition~\ref{prop:supp} hold for every field $\kappa$, and nothing in them
uses a normalization of $r$: the constant $M_r$ makes the argument uniform in $r\in(0,1)$. Only the final appeal to
Theorem~\ref{thm:FF} needs $\kappa$ finite. We do not need, and do not discuss, the analogue of
Theorem~\ref{thm:laurent} for infinite $\kappa$ with this particular pair $(E_1,B)$. For an infinite $\kappa$ of
characteristic $p$ the field $\kappa((X))$ is covered by Theorem~\ref{thm:main}, with the pair obtained by ascent from its
subfield $\F_p((X))$.
\end{remark}

\section{Ascent to an arbitrary base field}\label{sec:ascent}

\subsection{The Laurent subfield}
A metric space is \emph{spherically complete} if every nonempty family of closed balls, any two of which meet, has a
common point.

\begin{lemma}[discrete distances]\label{lem:discdist}
Let $Y$ be a complete ultrametric space whose distances lie in $r^\Z\cup\{0\}$ for some $r\in(0,1)$. Then $Y$ is
spherically complete.
\end{lemma}

\begin{proof}
We use two facts about closed balls $\cD(a,\rho)=\{y\in Y:\mathrm d(y,a)\le\rho\}$ in an ultrametric space: every point
of a ball is a center, and hence, if two closed balls meet, the one with the smaller radius is contained in the other.
Since all distances lie in $r^\Z\cup\{0\}$, we have $\cD(a,\rho)=\cD(a,\rho^\flat)$ with
$\rho^\flat:=\max\{s\in r^\Z\cup\{0\}:s\le\rho\}$, so we may assume that all radii lie in $r^\Z\cup\{0\}$. Let
$(\cD(a_i,\rho_i))_{i\in I}$ be a nonempty family of pairwise meeting closed balls. A subset of $r^\Z\cup\{0\}$ has a
least element unless its infimum is $0$ and not attained. \emph{If some $\rho_{i_0}$ is least,} then
$\cD(a_{i_0},\rho_{i_0})\subseteq\cD(a_i,\rho_i)$ for every $i$, and $a_{i_0}$ is a common point. \emph{Otherwise} all
$\rho_i>0$ and there are $i_1,i_2,\dots$ with $\rho_{i_n}$ strictly decreasing to $0$. The balls
$\cD(a_{i_n},\rho_{i_n})$ are nested, so $\mathrm d(a_{i_{n'}},a_{i_n})\le\rho_{i_n}$ for $n'\ge n$; the centers form a
Cauchy sequence with a limit $a$, which lies in every $\cD(a_{i_n},\rho_{i_n})$. For any $i$ choose $n$ with
$\rho_{i_n}\le\rho_i$; then $\cD(a_{i_n},\rho_{i_n})\subseteq\cD(a_i,\rho_i)$, so $a$ is a common point.
\end{proof}

\begin{remark}[the hypothesis (K2) of~\cite{M}]\label{rem:K2}
The construction of~\cite{M} assumes that the base field satisfies the hypothesis (K2) there: there are closed balls
$\cD(c_n,r_n)$ with $\abs{c_{n+1}-c_n}\le r_n$, $r_0>r_1>\cdots$ and $\bigcap_n\cD(c_n,r_n)=\emptyset$. An ultrametric
field (indeed any ultrametric space) satisfies (K2) if and only if it is not spherically complete. Indeed, the balls
in (K2) pairwise meet, since $\abs{c_n-c_m}\le\max_{m\le j<n}\abs{c_{j+1}-c_j}\le r_m$ for $m\le n$. Conversely, let a
family of pairwise meeting closed balls have no common point. A ball of least radius in the family would lie in all
the others (as in the proof of Lemma~\ref{lem:discdist}), and its center would be a common point; so the radii have
no least element, and every radius in the
family exceeds their infimum. A sequence of balls from the family whose radii decrease strictly to that infimum is
then nested, so it satisfies $\abs{c_{n+1}-c_n}\le r_n$, and it has empty intersection: a common point of the sequence
would lie in every ball of the family, because every ball of the family contains a ball of the sequence. This remark
is used only in \S\ref{sec:intro} and \S\ref{sec:lean}.
\end{remark}

\begin{lemma}[Laurent subfield]\label{lem:laurentsub}
Let $K'$ be a complete nontrivially normed field of characteristic $p$, and let $t\in K'$ with $0<\abs t=r<1$.
\begin{enumerate}
\item $K'$ is ultrametric, and $\abs c=1$ for every $c\in\F_p^\times$.
\item For $a_m,\dots,a_M\in\F_p$, not all zero, $\bigl\lvert\sum_{i=m}^Ma_it^i\bigr\rvert=r^{\min\{i:a_i\ne0\}}$.
\item Give $\F_p((X))$ the absolute value with $\abs X=r$. There is a continuous isometric ring homomorphism
$\mathrm{ev}_t\colon\F_p((X))\to K'$ such that
$\mathrm{ev}_t\bigl(\sum_{i\ge m}a_iX^i\bigr)=\sum_{i\ge m}a_it^i$. Its image
$K_1$ is a closed subfield of $K'$; it is the closure of $\F_p(t)$, and it is isometrically isomorphic to $\F_p((X))$.
\item $K_1$ is spherically complete.
\end{enumerate}
\end{lemma}

\begin{proof}
(1) $K'$ is ultrametric by Lemma~\ref{lem:ultra}, and $c^{p-1}=1$ gives $\abs c=1$.

(2) By (1) the nonzero terms have pairwise distinct absolute values $r^i$. In an ultrametric field $\abs{x+y}=\abs x$
whenever $\abs y<\abs x$, so the absolute value of the sum is the largest one, $r^{\min\{i:a_i\neq0\}}$.

(3) On Laurent polynomials, $\mathrm{ev}_t\colon\F_p[X,X^{-1}]\to K'$ is a ring homomorphism with
$\abs{\mathrm{ev}_t(g)}=r^{\ord g}=\abs g$, by (2) after multiplying by a power of $t$. Laurent polynomials are dense in
$\F_p((X))$ (truncations converge), and $K'$ is complete, so $\mathrm{ev}_t$ extends uniquely to a continuous map on
$\F_p((X))$. The extension is isometric, additive and multiplicative by continuity, and it maps the partial sums of
$\sum a_iX^i$ to those of $\sum a_it^i$. The image of the complete space $\F_p((X))$ under an isometry is complete, hence
closed; it is a subring, and a subfield because $\mathrm{ev}_t(g)\,\mathrm{ev}_t(g^{-1})=1$. It contains $\F_p(t)$, and
$\F_p[t,t^{-1}]\subseteq\F_p(t)$ is dense in it; so it is the closure of $\F_p(t)$.

(4) $K_1$ is complete by (3) and ultrametric by (1), and $\abs{K_1^\times}=r^\Z$, so its distances lie in
$r^\Z\cup\{0\}$. Apply Lemma~\ref{lem:discdist}.
\end{proof}

\subsection{Ingleton's theorem and its consequences}
Throughout the rest of this section we fix fields $K_1\subseteq K'$ with the following properties.
\begin{enumerate}
\item[(H1)] $K_1$, with the restricted absolute value, is nontrivially normed, complete and spherically complete.
\item[(H2)] $K'$ is ultrametric.
\end{enumerate}
By multiplication, $K'$ is an ultrametric normed $K_1$\nobreakdash-space. In our application $K'=\Kh$ and $K_1$ is the field of
Lemma~\ref{lem:laurentsub}.

\begin{lemma}\label{lem:findim}
Let $\mathbb K$ be a complete nontrivially normed field and $U$ a finite-dimensional normed $\mathbb K$\nobreakdash-space. Then every
linear functional on $U$ is bounded, and $U$ is complete.
\end{lemma}

\begin{proof}
By induction on $n=\dim U$; the case $n=0$ is trivial. Let $n\ge1$ and let $g\ne0$ be a linear functional on $U$. Its
kernel $H$ has dimension $n-1$, so it is complete by induction, hence closed in $U$. Choose $v$ with $g(v)=1$; then
$\delta:=\inf_{h\in H}\norm{v-h}>0$. Every $x\in U$ is $x=g(x)v+h$ with $h\in H$, and if $g(x)\ne0$ then
$\norm x=\abs{g(x)}\,\norm{v+g(x)^{-1}h}\ge\abs{g(x)}\delta$. So $\abs{g(x)}\le\delta^{-1}\norm x$. For completeness, the
coordinate map of a basis $U\to\mathbb K^n$ (max norm) is bounded by the first part, and its inverse
$(c_i)\mapsto\sum_ic_ie_i$ is bounded by $\sum_i\norm{e_i}$; so $U$ is bi-Lipschitz equivalent to the complete space
$\mathbb K^n$.
\end{proof}

\begin{theorem}[Ingleton~\cite{In}; see also~\cite{vR}]\label{thm:ingleton}
Let $K_1$ be a spherically complete ultrametric normed field, $V$ a $K_1$\nobreakdash-vector space with an ultrametric norm,
$V_0\subseteq V$ a subspace, $C\ge0$, and $T_0\colon V_0\to K_1$ linear with $\abs{T_0v}\le C\norm v$ for $v\in V_0$. Then
$T_0$ extends to a linear map $T\colon V\to K_1$ with $\abs{Tx}\le C\norm x$ for all $x\in V$.
\end{theorem}

We apply Theorem~\ref{thm:ingleton} only with the ultrametric domain $V=K'$ and the codomain $K_1$. It is never applied
to the spaces $E_1$, $F_1$ below, which need not be ultrametric.

\begin{corollary}[Ingleton projection]\label{cor:proj}
There is a $K_1$\nobreakdash-linear map $\varpi\colon K'\to K_1$ with $\varpi|_{K_1}=\id$ and $\abs{\varpi(\lambda)}\le\abs\lambda$ for all
$\lambda\in K'$.
\end{corollary}

\begin{proof}
Theorem~\ref{thm:ingleton} with $V=K'$, $V_0=K_1$, $T_0=\id$ and $C=1$; its hypotheses hold by (H1) and (H2).
\end{proof}

\begin{lemma}[extension]\label{lem:ext}
Let $\mathbb K\in\{K_1,K'\}$.
\begin{enumerate}
\item Let $V_0$ be a dense $\mathbb K$\nobreakdash-subspace of a normed $\mathbb K$\nobreakdash-space $V$, let $Y$ be a
$\mathbb K$\nobreakdash-Banach space and $T_0\colon V_0\to Y$ a $\mathbb K$\nobreakdash-linear map with $\norm{T_0x}\le C\norm x$. Then $T_0$ extends
uniquely to a continuous map $T\colon V\to Y$, which is $\mathbb K$\nobreakdash-linear with $\norm{Tx}\le C\norm x$.
\item Let $V_1,\dots,V_k$ be normed $K'$\nobreakdash-spaces with dense subspaces $V_i^0$, $Y$ a $K'$\nobreakdash-Banach space, and
$M\colon\prod_iV_i^0\to Y$ a $K'$\nobreakdash-multilinear map with $\norm{M(d)}\le C\prod_i\norm{d_i}$. Then $M$ extends uniquely to a
continuous $K'$\nobreakdash-multilinear map $\overline M\colon\prod_iV_i\to Y$ with $\norm{\overline M}\le C$.
\end{enumerate}
\end{lemma}

\begin{proof}
(1) $T_0$ is uniformly continuous and $Y$ is complete, so it extends uniquely to a continuous map; linearity and the
bound pass to limits by continuity of addition and of scalar multiplication. (2) For $d,d'\in\prod_iV_i^0$ with all
$\norm{d_i},\norm{d'_i}\le\rho$ the telescoping estimate
$\norm{M(d)-M(d')}\le C\rho^{k-1}\sum_i\norm{d_i-d'_i}$ holds. So $M$ is uniformly continuous on bounded subsets of
$\prod_iV_i^0$, and it extends uniquely to a continuous map on their closures, which cover $\prod_iV_i$; multilinearity
and the bound pass to the limit.
\end{proof}

\subsection{Projective base change}
For a normed $K_1$\nobreakdash-space $V$ let $V_{K'}:=V\otimes_{K_1}K'$, a $K'$\nobreakdash-vector space through the second factor, and put
\[
\norm u_\pi:=\inf\Bigl\{\sum_j\norm{x_j}\,\abs{\lambda_j}\ :\ u=\sum_jx_j\otimes\lambda_j\Bigr\},
\]
the infimum over finite decompositions, with an ordinary sum. For a bounded $K_1$\nobreakdash-linear $\psi\colon K'\to K_1$ let
$T_\psi\colon V_{K'}\to V$ be the $K_1$\nobreakdash-linear map with $T_\psi(x\otimes\lambda)=\psi(\lambda)x$; it exists because
$(x,\lambda)\mapsto\psi(\lambda)x$ is $K_1$\nobreakdash-bilinear. Let $\iota_V\colon V\to V_{K'}$, $x\mapsto x\otimes1$; it is
$K_1$\nobreakdash-linear, since $cx\otimes1=x\otimes c=c\cdot(x\otimes1)$ for $c\in K_1$.

\begin{lemma}[projective base change]\label{lem:bc}
Let $V$ be a normed $K_1$\nobreakdash-space; no ultrametric hypothesis is made on $V$.
\begin{enumerate}
\item $\norm{T_\psi u}\le\norm\psi\,\norm u_\pi$ for every bounded $K_1$\nobreakdash-linear $\psi\colon K'\to K_1$.
\item $\norm\cdot_\pi$ is a $K'$\nobreakdash-norm on $V_{K'}$.
\item $\iota_V$ is isometric.
\item Let $V\widehat\otimes_\pi K'$ be the completion of $(V_{K'},\norm\cdot_\pi)$. If $V$ is complete, then $T_\varpi$, with
$\varpi$ as in Corollary~\ref{cor:proj}, extends to a $K_1$\nobreakdash-linear map $\Pi_V\colon V\widehat\otimes_\pi K'\to V$ with
$\norm{\Pi_V}\le1$ and $\Pi_V\circ\iota_V=\id_V$.
\end{enumerate}
\end{lemma}

\begin{proof}
(1) For any decomposition $u=\sum_jx_j\otimes\lambda_j$ we have
\[
\norm{T_\psi u}=\Bigl\lVert\sum_j\psi(\lambda_j)x_j\Bigr\rVert\le\sum_j\norm\psi\abs{\lambda_j}\norm{x_j};
\]
take the infimum.

(2) Multiplying the second factors by $\mu\in K'$ gives $\norm{\mu u}_\pi\le\abs\mu\norm u_\pi$, and applying this to
$\mu^{-1}$ gives equality for $\mu\ne0$; concatenating decompositions gives the triangle inequality. For positivity let
$u\neq0$ and write $u=\sum_{j=1}^nx_j\otimes\lambda_j$ with $\lambda_1,\dots,\lambda_n$ linearly independent over $K_1$
(rewrite any representation in a basis of the $K_1$\nobreakdash-span of its second factors). Since $u\ne0$, some $x_{j_0}\ne0$. The
$K_1$\nobreakdash-span $U$ of the $\lambda_j$ is a finite-dimensional normed $K_1$\nobreakdash-space, so the coordinate functional
$\psi_0\colon U\to K_1$, $\psi_0(\lambda_j)=\delta_{jj_0}$, is bounded by some $C$ (Lemma~\ref{lem:findim}; $K_1$ is complete
by (H1)). By Theorem~\ref{thm:ingleton}, applied to the ultrametric $K_1$\nobreakdash-space $K'$, it extends to $\psi\colon K'\to K_1$
with $\norm\psi\le C$. Then $T_\psi u=x_{j_0}$, and (1) gives $0<\norm{x_{j_0}}\le C\norm u_\pi$. No functional on $V$ is
used.

(3) $\norm{x\otimes1}_\pi\le\norm x$, and $\norm x=\norm{T_\varpi(x\otimes1)}\le\norm{x\otimes1}_\pi$ by (1), since
$\norm\varpi\le1$ and $\varpi(1)=1$.

(4) By (1), $T_\varpi$ is $K_1$\nobreakdash-linear and $1$\nobreakdash-Lipschitz on the dense subspace $V_{K'}$, with values in the complete space
$V$; apply Lemma~\ref{lem:ext}(1) with $\mathbb K=K_1$. Finally $\Pi_V\iota_Vx=T_\varpi(x\otimes1)=\varpi(1)x=x$.
\end{proof}

\emph{The data.} Let $E_1$ be a normed $K_1$\nobreakdash-space and $F_1$ a $K_1$\nobreakdash-\emph{Banach} space. Put
\[
E:=E_1\widehat\otimes_\pi K',\qquad F:=F_1\widehat\otimes_\pi K',
\]
the completions of $((E_1)_{K'},\norm\cdot_\pi)$ and $((F_1)_{K'},\norm\cdot_\pi)$. They are $K'$\nobreakdash-Banach spaces;
$\iota_E\colon E_1\to E$ and $\iota_F\colon F_1\to F$ are $K_1$\nobreakdash-linear isometries, and $\Pi_F\colon F\to F_1$ is $K_1$\nobreakdash-linear
with $\norm{\Pi_F}\le1$ and $\Pi_F\circ\iota_F=\id_{F_1}$ (Lemma~\ref{lem:bc}).

\begin{lemma}[operators]\label{lem:ops}
For $f\in\cL_{K_1}(E_1,E_1)$ the map $f\otimes\id$ on $(E_1)_{K'}$ is $K'$\nobreakdash-linear with
$\norm{(f\otimes\id)u}_\pi\le\norm f\norm u_\pi$. Its extension $f_{K'}\in\cL_{K'}(E,E)$ satisfies $\norm{f_{K'}}\le\norm f$
and $f_{K'}\circ\iota_E=\iota_E\circ f$, and $f\mapsto f_{K'}$ is $K_1$\nobreakdash-linear.
\end{lemma}

\begin{proof}
Apply $f$ to the first factors of any decomposition, and extend by Lemma~\ref{lem:ext}(1) with $\mathbb K=K'$, since $E$ is
complete. The identity $f_{K'}\iota_E=\iota_Ef$ and the $K_1$\nobreakdash-linearity in $f$ are visible on elementary tensors.
\end{proof}

\begin{lemma}[forms]\label{lem:forms-bc}
For $m\in\Alt^k_{K_1}(E_1;F_1)$ there is $m_{K'}\in\Alt^k_{K'}(E;F)$ with $\norm{m_{K'}}\le\norm m$ and
\[
m_{K'}(\iota_Ex_1,\dots,\iota_Ex_k)=\iota_F\bigl(m(x_1,\dots,x_k)\bigr)\qquad(x_i\in E_1),
\]
and $m\mapsto m_{K'}$ is $K_1$\nobreakdash-linear.
\end{lemma}

\begin{proof}
Fix a $K_1$\nobreakdash-basis $(\lambda_\beta)_{\beta\in\mathfrak B}$ of $K'$. Every $u\in(E_1)_{K'}$ has a unique expansion
$u=\sum_\beta u^\beta\otimes\lambda_\beta$ with $u^\beta\in E_1$, finitely many nonzero, and $u\mapsto u^\beta$ is
$K_1$\nobreakdash-linear. Define
\[
m^{\mathrm{alg}}(u_1,\dots,u_k):=\sum_{\beta_1,\dots,\beta_k}m\bigl(u_1^{\beta_1},\dots,u_k^{\beta_k}\bigr)\otimes\lambda_{\beta_1}\cdots\lambda_{\beta_k}\in(F_1)_{K'} .
\]
\textup{(i)} $m^{\mathrm{alg}}$ is $K_1$\nobreakdash-multilinear.

\textup{(ii)} $m^{\mathrm{alg}}(x_1\otimes\mu_1,\dots,x_k\otimes\mu_k)=m(x_1,\dots,x_k)\otimes\mu_1\cdots\mu_k$ for $x_i\in E_1$ and
$\mu_i\in K'$. Indeed, write $\mu_i=\sum_\beta c_{i\beta}\lambda_\beta$ with $c_{i\beta}\in K_1$; then
$x_i\otimes\mu_i=\sum_\beta(c_{i\beta}x_i)\otimes\lambda_\beta$ and
$m^{\mathrm{alg}}=\sum_\beta\prod_ic_{i\beta_i}\,m(x)\otimes\prod_i\lambda_{\beta_i}=m(x)\otimes\prod_i\bigl(\sum_\beta c_{i\beta}\lambda_\beta\bigr)$.

\textup{(iii)} $m^{\mathrm{alg}}$ is $K'$\nobreakdash-multilinear. By \textup{(i)} and additivity it suffices to check $K'$\nobreakdash-homogeneity in
one slot when every slot is an elementary tensor, and there \textup{(ii)} gives it.

\textup{(iv)} $\norm{m^{\mathrm{alg}}(u_1,\dots,u_k)}_\pi\le\norm m\prod_i\norm{u_i}_\pi$. For arbitrary decompositions
$u_i=\sum_jx_{ij}\otimes\mu_{ij}$, multilinearity and \textup{(ii)} give
$m^{\mathrm{alg}}(u)=\sum_{j_1,\dots,j_k}m(x_{1j_1},\dots,x_{kj_k})\otimes\prod_i\mu_{ij_i}$, of projective norm at most
$\norm m\prod_i\bigl(\sum_j\norm{x_{ij}}\abs{\mu_{ij}}\bigr)$; take the infimum in each slot.

\textup{(v)} $m^{\mathrm{alg}}$ is alternating in the strong sense. Let slots $a\ne b$ both carry
$u=\sum_jx_j\otimes\mu_j$, and expand every slot into elementary tensors. By \textup{(ii)} each term is
$m(\dots)\otimes(\text{scalar})$. Terms with the same index $j$ in slots $a$ and $b$ contain $m$ with a repeated
argument and vanish. The terms with indices $(j,j')$, $j\ne j'$, in slots $(a,b)$ pair off with those with indices
$(j',j)$: the scalars agree, and the values of $m$ are negatives of each other because $m$ is antisymmetric. So the total
is $0$. This argument is valid in every characteristic, including $2$; checking only the full diagonal
$m^{\mathrm{alg}}(u,\dots,u)=0$ would not suffice for $k\ge3$.

\textup{(vi)} By \textup{(iii)}, \textup{(iv)} and Lemma~\ref{lem:ext}(2) ($F$ is complete), $m^{\mathrm{alg}}$ extends to a continuous
$K'$\nobreakdash-multilinear map $m_{K'}\colon E^k\to F$ with $\norm{m_{K'}}\le\norm m$. It is alternating: if $v\in E^k$ has
$v_a=v_b$, approximate $v_a=v_b$ by a single sequence in $(E_1)_{K'}$ and the other $v_i$ by sequences in $(E_1)_{K'}$;
$m^{\mathrm{alg}}$ vanishes along the approximants by \textup{(v)}, so $m_{K'}(v)=0$ by continuity.

\textup{(vii)} $m_{K'}(\iota_Ex_1,\dots,\iota_Ex_k)=m^{\mathrm{alg}}(x_1\otimes1,\dots,x_k\otimes1)=m(x)\otimes1=\iota_F(m(x))$ by
\textup{(ii)}. $K_1$\nobreakdash-linearity of $m\mapsto m^{\mathrm{alg}}$ is visible in the formula, and it passes to the extensions by
uniqueness.
\end{proof}

\begin{proposition}[descent of a lift]\label{prop:ascent}
Assume \textup{(H1)} and \textup{(H2)}. Let $E_1$ be a normed $K_1$\nobreakdash-space, $F_1$ a $K_1$\nobreakdash-Banach space, and $E,F$ as
above. If $Q^{K'}_{E,E,F}$ admits a bounded $k$\nobreakdash-linear lift $P$ over $K'$, then $Q^{K_1}_{E_1,E_1,F_1}$ admits a
bounded $k$\nobreakdash-linear lift $P_1$ over $K_1$ with $\norm{P_1}\le\norm P$. Equivalently, by Proposition~\ref{prop:lift} over $K_1$ and over
$K'$: if $Q^{K_1}_{E_1,E_1,F_1}$ is analytic at no point, then $Q^{K'}_{E,E,F}$ is analytic at no point either.
\end{proposition}

\begin{proof}
For $f_1,\dots,f_k\in\cL_{K_1}(E_1,E_1)$, $m\in\Alt^k_{K_1}(E_1;F_1)$ and $x\in E_1^k$ define
\[
P_1(f_1,\dots,f_k)(m)(x_1,\dots,x_k):=\Pi_F\Bigl(P\bigl((f_1)_{K'},\dots,(f_k)_{K'}\bigr)(m_{K'})(\iota_Ex_1,\dots,\iota_Ex_k)\Bigr).
\]
\emph{It is well defined.} $P((f_i)_{K'})(m_{K'})\in\Alt^k_{K'}(E;F)$ is $K'$\nobreakdash-multilinear, hence $K_1$\nobreakdash-multilinear, and
$\iota_E$ and $\Pi_F$ are $K_1$\nobreakdash-linear. So $x\mapsto P_1(f)(m)(x)$ is $K_1$\nobreakdash-multilinear, vanishes when two $x_i$ coincide
(then the $\iota_Ex_i$ coincide), and has norm at most
$\norm{\Pi_F}\norm P\prod_i\norm{(f_i)_{K'}}\norm{m_{K'}}\le\norm P\prod_i\norm{f_i}\norm m$ by
Lemmas~\ref{lem:bc}--\ref{lem:forms-bc}. So $P_1(f)(m)\in\Alt^k_{K_1}(E_1;F_1)$. It is $K_1$\nobreakdash-linear in $m$, because
$m\mapsto m_{K'}$ is $K_1$\nobreakdash-linear and $P(\dots)$ is $K'$\nobreakdash-linear, and $K_1$\nobreakdash-multilinear in $(f_1,\dots,f_k)$, because
$f\mapsto f_{K'}$ is $K_1$\nobreakdash-linear and $P$ is $K'$\nobreakdash-multilinear. Hence $P_1$ is a bounded $k$\nobreakdash-linear map over $K_1$ with
$\norm{P_1}\le\norm P$.

\emph{The diagonal.} Using $P(g,\dots,g)=Q^{K'}(g)$ with $g=f_{K'}$, Lemma~\ref{lem:ops} and Lemma~\ref{lem:forms-bc},
\begin{align*}
P_1(f,\dots,f)(m)(x)&=\Pi_F\bigl(m_{K'}(f_{K'}\iota_Ex_1,\dots,f_{K'}\iota_Ex_k)\bigr)
=\Pi_F\bigl(m_{K'}(\iota_Efx_1,\dots,\iota_Efx_k)\bigr)\\
&=\Pi_F\bigl(\iota_F(m(fx_1,\dots,fx_k))\bigr)=m(fx_1,\dots,fx_k)=Q^{K_1}(f)(m)(x).\qedhere
\end{align*}
\end{proof}

\begin{remark}\label{rem:ascent}
Proposition~\ref{prop:ascent} is characteristic-free and uses no hypothesis on $k$. It transports negative results
upward only, from $K_1$ to $K'$. Ingleton's theorem is used only on $K'$, so $F_1$ may be, and in our application is,
non-ultrametric. The completeness of $F_1$ is used essentially: it is what makes the retraction $\Pi_F$ land in $F_1$.
\end{remark}

\subsection{Incomplete base fields}

\begin{lemma}\label{lem:descent}
Let $K$ be a nontrivially normed field with completion $\Kh$, and let $E,E',F$ be $\Kh$\nobreakdash-Banach spaces, regarded as
$K$\nobreakdash-Banach spaces by restriction of scalars (with the same norms). Then $\cL_K(E,E')=\cL_{\Kh}(E,E')$,
$\Alt^k_K(E;F)=\Alt^k_{\Kh}(E;F)$ and $\Alt^k_K(E';F)=\Alt^k_{\Kh}(E';F)$ with the same norms, $Q^K_{E,E',F}=Q^{\Kh}_{E,E',F}$,
and every bounded $k$\nobreakdash-linear lift of $Q^K_{E,E',F}$ over $K$ is a bounded $k$\nobreakdash-linear lift of $Q^{\Kh}_{E,E',F}$ over
$\Kh$.
\end{lemma}

\begin{proof}
A continuous $K$\nobreakdash-linear map $g$ between $\Kh$\nobreakdash-normed spaces is $\Kh$\nobreakdash-linear: for $\lambda\in\Kh$ choose $\lambda_n\in K$
with $\lambda_n\to\lambda$; then $\lambda_nx\to\lambda x$ because $\norm{\lambda_nx-\lambda x}=\abs{\lambda_n-\lambda}\norm x$,
so $g(\lambda x)=\lim g(\lambda_nx)=\lim\lambda_ng(x)=\lambda g(x)$. Likewise a continuous $K$\nobreakdash-multilinear map is
$\Kh$\nobreakdash-multilinear, slot by slot. So the spaces of bounded maps coincide as sets, with the same sup-formula norms, and
they are $\Kh$\nobreakdash-Banach spaces. The map $Q$ is given by the same formula over both fields. A bounded $k$\nobreakdash-linear map
$P$ over $K$ between these spaces is continuous in each slot, hence $k$\nobreakdash-linear over $\Kh$.
\end{proof}

\needspace{8\baselineskip}
\section{Proofs of the Main Theorem and the Corollary}\label{sec:proofs}
We first show that the target has no equivalent ultrametric norm, and then assemble the proofs of
Theorem~\ref{thm:main} and Corollary~\ref{cor:class}.

\subsection{The target has no equivalent ultrametric norm}

\begin{proposition}\label{prop:nonultra}
Let $\kappa$ be any field, $k\ge2$, $r\in(0,1)$, and let $K_1=\kappa((X))$, $E_1$, $E_0$ and $B$ be as in
\S\ref{sec:laurent}. Let $e_i\in E_0$ be the unit vectors and $\omega_j:=e_{kj+1}\wedge\dots\wedge e_{kj+k}\in B$
($j\ge0$), wedges on pairwise disjoint blocks of indices. Then $\norm{\omega_j}_B=1$ and
\[
\frac N{M_r}\ \le\ \Bigl\lVert\sum_{j<N}\omega_j\Bigr\rVert_B\ \le\ N\qquad(N\ge1).
\]
Hence the norm of $B$ is not ultrametric, and $B$ admits no equivalent ultrametric norm. The same holds for
$F=B\widehat\otimes_\pi K'$ as in \S\ref{sec:ascent}, for any $K'\supseteq K_1$ satisfying \textup{(H1)} and \textup{(H2)},
regarded as a normed space over $K'$ or over any subfield of $K'$.
\end{proposition}

\begin{proof}
The upper bounds follow from the defining decompositions, since unit vectors have norm $1$: each $\omega_j$ is one
wedge of unit vectors, and $\sum_{j<N}\omega_j$ is a sum of $N$ of them. Also
$\norm{\omega_j}_B\ge\norm{J\omega_j}_\infty\ge1$, because the entry of $J\omega_j$ at the block indices is $\det I=1$.

Put $\omega:=\sum_{j<N}\omega_j$. It is the image of the element $\sum_{j<N}\omega_j$ of $\Lambda^k_\kappa E_0$, again
denoted $\omega$; its array $J\omega=\Omega^\kappa(\omega)$ has entries in $\kappa$, so $\coeff_0$ fixes it and
$\eta(\omega)=\omega$ by uniqueness in Proposition~\ref{prop:supp}. We bound $\sdim(\omega)$ from below with Lemma~\ref{lem:contr}. Fix a
block $j<N$ and an index $i_0$ in it, and let $\varphi$ consist of the $k-1$ coordinate functionals $e^*_{i'}$ of the
other indices $i'$ of block $j$. For a block $j'\ne j$ all $\varphi_a$ vanish on the vectors of $\omega_{j'}$, so every
$(k-1)\times(k-1)$ matrix in the formula for $c_\varphi(\omega_{j'})$ is zero and $c_\varphi(\omega_{j'})=0$ (here $k\ge2$
is used). In $c_\varphi(\omega_j)$, the term of the vector $e_{i_0}$ has a permutation matrix, with determinant $\pm1$, and
every other term has a matrix containing the column of $e_{i_0}$, on which all $\varphi_a$ vanish. So
$c_\varphi(\omega)=\pm e_{i_0}$. The span of the $c_\varphi(\omega)$ therefore contains the $kN$ unit vectors of the first $N$
blocks, and $\sdim(\omega)\ge kN$. Now~\eqref{eq:suppest} gives
$kN\le\sdim(\eta(\omega))\le kM_r\norm\omega_B$.

If the norm of $B$ were ultrametric, then $\norm\omega_B\le\max_j\norm{\omega_j}_B=1<N/M_r$ for $N>M_r$. If $\norm\cdot_u$ were
an equivalent ultrametric norm, with $\norm\cdot_B\le c\norm\cdot_u\le c'\norm\cdot_B$, then
$\norm{\sum_{j<N}\omega_j}_B\le c\,\max_j\norm{\omega_j}_u\le c'$ for every $N$, contradicting $N/M_r\to\infty$. (Only
the ultrametric inequality of $\norm\cdot_u$ is used, not its homogeneity.)

For $F$: $\iota_F\colon B\to F$ is a $K_1$\nobreakdash-linear isometry (Lemma~\ref{lem:bc}(3)), so the vectors $\iota_F\omega_j$
satisfy $\norm{\iota_F\omega_j}=1$ and $\norm{\sum_{j<N}\iota_F\omega_j}\ge N/M_r$, and the same argument applies.
Regarding $F$ over a subfield does not change its norm.
\end{proof}

If $r\le\frac12$, then $M_r=1$ and $\norm{\sum_{j<N}\omega_j}_B=N$ exactly.

\subsection{Proof of Theorem~\ref{thm:main}}
Let $K$ be a nontrivially normed field of characteristic $p>0$, and $k\ge p$.
\begin{enumerate}
\item \emph{The base.} $K$ is ultrametric (Lemma~\ref{lem:ultra}), and its completion $\Kh$ is a complete, nontrivially
normed, ultrametric field of characteristic $p$. Choose $t\in K$ with $0<\abs t=r<1$; it exists because $K$ is
nontrivially normed.
\item \emph{The Laurent subfield.} Let $K_1\subseteq\Kh$ be the closure of $\F_p(t)$. By Lemma~\ref{lem:laurentsub} with
$K'=\Kh$, it is a closed subfield, isometrically isomorphic to $\F_p((X))$ with $\abs X=r$ via $\mathrm{ev}_t$, and it is
complete and spherically complete.
\item \emph{No lift over $K_1$.} Put $E_1:=\linf(\N,K_1)$, and let $B$ be the completion of $\Lambda^k_{K_1}E_1$ under the
projective exterior norm. Theorem~\ref{thm:laurent} with $\kappa=\F_p$ holds for $\F_p((X))$, and it transports along the
isometric isomorphism $\mathrm{ev}_t$: via $\mathrm{ev}_t^{-1}$ a normed $\F_p((X))$\nobreakdash-space is a normed $K_1$\nobreakdash-space with the
same group and norm, $\linf(\N,\F_p((X)))\cong\linf(\N,K_1)$ isometrically, and this identification preserves exterior
powers, projective norms, the completion $B$, the spaces $\cL$ and $\Alt^k$, the maps $Q$ and their bounded lifts. So
$Q^{K_1}_{E_1,E_1,B}$ has no bounded $k$\nobreakdash-linear lift over $K_1$.
\item \emph{Ascent to $\Kh$.} Hypothesis (H1) holds for $K_1\subseteq\Kh$ by (2), since $\abs t=r\in(0,1)$, and (H2)
holds by (1). The space $F_1:=B$ is complete. By Proposition~\ref{prop:ascent} with $K'=\Kh$, the map
$Q^{\Kh}_{E,E,F}$ for
\[
E:=E_1\widehat\otimes_\pi\Kh,\qquad F:=B\widehat\otimes_\pi\Kh
\]
has no bounded $k$\nobreakdash-linear lift over $\Kh$.
\item \emph{Descent to $K$.} Regard $E$ and $F$ as $K$\nobreakdash-Banach spaces. By Lemma~\ref{lem:descent}, a bounded
$k$\nobreakdash-linear lift of $Q^K_{E,E,F}$ over $K$ would be a bounded $k$\nobreakdash-linear lift of $Q^{\Kh}_{E,E,F}$ over $\Kh$. So
there is none.
\item \emph{Nowhere analytic.} By Proposition~\ref{prop:lift}, which holds over the possibly incomplete field $K$,
$Q^K_{E,E,F}$ is analytic at no point, and \texttt{ContDiffAt K $\omega$} fails at every point. For a finite set $\iota$
with $\abs\iota=k$ the same holds, by the identification of \S\ref{sec:prelim}.
\item \emph{The target.} Since $k\ge p\ge2$, Proposition~\ref{prop:nonultra} (with $\kappa=\F_p$, transported along $\mathrm{ev}_t$ as in
step~(3), and $K'=\Kh$, regarding $F$ over $K$) shows that $F$ has no equivalent ultrametric norm.\qed
\end{enumerate}

\begin{remark}\label{rem:spaces}
If $K_1=\Kh$, for example if $K=\F_p((u))$ and $t=u$, then $E=E_1$ and $F=B$ (Lemma~\ref{lem:bc}(3) and $x\otimes\lambda=\lambda
x\otimes1$), so $E$ is ultrametric. For $K=\F_q((u))$, $q=p^n$ with the $u$\nobreakdash-adic absolute value, one can also bypass the
ascent: Theorem~\ref{thm:laurent} with $\kappa=\F_q$ and $X=u$ gives $E=\linf(\N,K)$ and $F=B$ directly. We do not claim
that $E$ is ultrametric in general; the proof never needs it to be.
\end{remark}

\subsection{Proof of Corollary~\ref{cor:class}}
Put $k:=\abs\iota$.

(1) Assume $k!\ne0$ in $K$, and let $E,E',F$ be normed. Define
\[
P(f_1,\dots,f_k)(m)(x_1,\dots,x_k):=\frac1{k!}\sum_{\sigma\in S_k}m\bigl(f_{\sigma(1)}x_1,\dots,f_{\sigma(k)}x_k\bigr).
\]
For fixed $f$ and $m$ this is $k$\nobreakdash-linear in $x$; each of the $k!$ terms of the sum has norm at most
$\norm m\prod_i\norm{f_i}\prod_j\norm{x_j}$, so it is bounded.
It is alternating: if $x_a=x_b=y$ with $a\ne b$, let $\tau$ be the transposition of $a$ and $b$; the terms for $\sigma$ and
$\sigma\circ\tau$ are $m(\dots,f_{\sigma(a)}y,\dots,f_{\sigma(b)}y,\dots)$ and $m(\dots,f_{\sigma(b)}y,\dots,f_{\sigma(a)}y,\dots)$
(slots $a$ and $b$ shown), which cancel because $m$ is antisymmetric, and $\sigma\mapsto\sigma\circ\tau$ is a
fixed-point-free involution of $S_k$. So $P(f)(m)\in\Alt^k(E;F)$; it is linear in $m$, and $P$ is a bounded $k$\nobreakdash-linear
map. On the diagonal all $k!$ terms equal $m(fx_1,\dots,fx_k)$, so $P(f,\dots,f)=Q(f)$. By
Proposition~\ref{prop:lift}(1), $Q$ is given at every point by a finite power series; for $\iota=\emptyset$, where
Proposition~\ref{prop:lift} assumes $k\ge1$, this is trivial, because $Q$ is then constant. No completeness is used.

(2) Assume $k!=0$ in $K$. In characteristic $0$ the positive integer $k!$ is nonzero in $K$, so $\operatorname{char}K=p>0$,
and $k!=0$ in $\F_p$ means $p\le k$. Theorem~\ref{thm:main} gives the required spaces.

Together, (1) and (2) give the equivalence.\qed

\needspace{8\baselineskip}
\section{Complements and open problems}\label{sec:complements}

\subsection*{Theorem A and targets over discretely valued bases}
The following positive result is the third item of \S\ref{sec:intro}. A special case is stated without proof
in~\cite[Remark~16]{M}. We quote the general statement from an unpublished Lean formalization (\S\ref{sec:lean}) and
do not prove it. It is used only in this subsection, and not in the proofs of Theorem~\ref{thm:main} and Corollary~\ref{cor:class}.

\begin{thmA}[{\cite[Remark~16]{M}} for a special case; formalized, \S\ref{sec:lean}]
Let $K$ be an ultrametric nontrivially normed field, let $E,E'$ be normed $K$\nobreakdash-spaces, and let $F$ be a normed $K$\nobreakdash-space
whose norm is ultrametric and which is spherically complete as a metric space. Then $Q^K_{E,E',F}$ is $C^n$ for every
$n\in\N\cup\{\infty,\omega\}$; in particular it is analytic at every point.
\end{thmA}

\begin{proposition}[targets over discretely valued bases]\label{prop:discrete}
Assume Theorem~A. Let $K$ be an ultrametric normed field with $\abs{K^\times}=r^\Z$ for some $r\in(0,1)$, complete or
not; for instance $K=\F_q((t))$ or $K=\kappa((t))$. Let $E,E'$ be normed $K$\nobreakdash-spaces and $F$ a $K$\nobreakdash-Banach space that
admits an equivalent ultrametric norm. Then $Q^K_{E,E',F}$ is analytic at every point.
\end{proposition}

\begin{proof}
\emph{Rounding.} Let $\norm\cdot_u$ be an ultrametric norm on $F$ equivalent to its given norm. For $y\ne0$ let
$\norm y_u^\sharp$ be the least element of $\{s\in r^\Z:s\ge\norm y_u\}$, which exists because $\norm y_u>0$, and put
$\norm0_u^\sharp:=0$. Then $\norm y_u\le\norm y_u^\sharp\le r^{-1}\norm y_u$. Rounding up to $r^\Z$ is monotone and
commutes with multiplication by elements of $r^\Z$. Since $\abs\lambda\in r^\Z$ for $\lambda\in K^\times$, this gives
$\norm{\lambda y}_u^\sharp=\abs\lambda\norm y_u^\sharp$, and
$\norm{y+y'}_u^\sharp\le\max(\norm y_u^\sharp,\norm{y'}_u^\sharp)$ follows from the same inequality for
$\norm\cdot_u$. So $\norm\cdot_u^\sharp$ is an ultrametric norm on $F$, equivalent to the given one. Hence
$F':=(F,\norm\cdot_u^\sharp)$ is complete, and its distances lie in $r^\Z\cup\{0\}$.

\emph{Theorem A.} By Lemma~\ref{lem:discdist}, $F'$ is spherically complete, so $Q':=Q^K_{E,E',F'}$ is analytic at
every point by Theorem~A.

\emph{Equivalent norms.} Boundedness and alternation of a multilinear map do not depend on the choice between
equivalent norms on the target. So $\Alt^k(E';F)=\Alt^k(E';F')$ and $\Alt^k(E;F)=\Alt^k(E;F')$ as vector spaces, with
equivalent operator norms. Hence $Z:=\cL(\Alt^k(E';F),\Alt^k(E;F))$ and $Z':=\cL(\Alt^k(E';F'),\Alt^k(E;F'))$ are the
same vector space with equivalent norms, and $Q=Q'$ as maps from $\cL(E,E')$ to it. A power series of $Q'$ at a point,
with radius $R$ and $\sup_n\norm{p_n}R^n<\infty$ for the norm of $Z'$, is a power series of $Q$ at that point with
the same radius: the operator norms of the $p_n$ change by at most a constant factor, and convergence in $Z'$ is
convergence in $Z$.
\end{proof}

Consequently, over such $K$, the target $F$ of every Banach counterexample $(E,E',F)$ admits no equivalent ultrametric
norm, as the target of Theorem~\ref{thm:main} does (Proposition~\ref{prop:nonultra}).

\begin{remark}[necessary, but not sufficient]\label{rem:notsuff}
Over discretely valued bases, having no equivalent ultrametric norm is therefore necessary for the target of a
Banach counterexample. The completeness of $F$ is used in Proposition~\ref{prop:discrete} through
Lemma~\ref{lem:discdist}, and the statement concerns complete targets only. The condition is not sufficient. For $(p,k)=(5,3)$, the space $B$ of \S\ref{ssec:projnorm}, built in degree~$3$
over $K_1=\F_5((X))$, has no equivalent ultrametric norm by Proposition~\ref{prop:nonultra}, which needs only
$k\ge2$. Yet $Q^{K_1}_{E_1,E_1,B}$ is analytic by Corollary~\ref{cor:class}(1), since $3!=1$ in $\F_5$. The
obstruction is the vanishing of $k!$. Indeed, feeding the bounded lift of Corollary~\ref{cor:class}(1) into
\S\ref{ssec:Psi} with $\kappa=\F_5$ produces a map $\Psi$ with the properties of Lemma~\ref{lem:Psi}, which hold in
every degree. So $\Psi$ satisfies all the hypotheses of Theorem~\ref{thm:FF}(2) with $L=\F_5$ except $k!=0$; compare
Remark~\ref{rem:FFunit}.
\end{remark}

\begin{remark}[an explicit lift]\label{rem:sorted}
Let $K$ be ultrametric, let $E$ have a basis $(e_i)_{i\in I}$ that is orthogonal for the norm (every $x$ is
$\sum_ie^*_i(x)e_i$, and $\norm{\sum_ic_ie_i}=\sup_i\abs{c_i}\norm{e_i}$), and let $F$ be complete and ultrametric. Fix a
total order on $I$ and put
\[
P(f_1,\dots,f_k)(m)(x^1,\dots,x^k):=\sum_{s_1<\dots<s_k}\det\bigl(e^*_{s_a}(x^b)\bigr)_{a,b}\;m\bigl(f_1e_{s_1},\dots,f_ke_{s_k}\bigr).
\]
For every $\varepsilon>0$ only finitely many terms have norm at least $\varepsilon$, so the family is summable in the
complete ultrametric $F$. By orthogonality $\abs{e^*_s(x)}\norm{e_s}\le\norm x$, so by the ultrametric
inequality in $K$ every term has norm at most $\norm m\prod_a\norm{f_a}\prod_b\norm{x^b}$, and $\norm P\le1$ by the
ultrametric inequality in $F$. The output is alternating because of the determinant; and on the diagonal the expansion
$m(fx^1,\dots,fx^k)=\sum_{s_1<\dots<s_k}\det(e^*_{s_a}(x^b))\,m(fe_{s_1},\dots,fe_{s_k})$ of an alternating map shows
$P(f,\dots,f)=Q(f)$. This \emph{sorted lift} is a bounded lift in every characteristic and degree, with no division by
$k!$ and no spherical completeness; we do not use it. Over a complete discretely valued field every ultrametric Banach
space is topologically isomorphic to a space $c_0(I)$ of null families~\cite[\S1, Prop.~1 and its Corollaire]{Se}, so
it has such a basis for an equivalent norm.
\end{remark}

\begin{remark}[ultrametric targets over bases that are not spherically complete; announced]\label{rem:R}
We state the following without proof. It is not proved in this paper, and nothing here depends on it. Let $k\ge2$, and
let $K$ be a complete nontrivially normed field with $k!=0$ in $K$ (so $K$ is ultrametric, by Lemma~\ref{lem:ultra})
that is not spherically complete. For a cardinal $\mu$ put $\beth_0(\mu):=\mu$ and
$\beth_{n+1}(\mu):=2^{\beth_n(\mu)}$, and let $\mu^+$ be the successor cardinal of $\mu$. Let $I$ be a set of
cardinality at least $\beth_{3k-1}(\abs K)^+$, put $E:=\linf(I,K)$, and let $F_I$ be the closure in $\linf(I^k,K)$ of
the span of the determinant arrays $(i_1,\dots,i_k)\mapsto\det(x_b(i_a))_{a,b}$, $x\in E^k$. Then $E$ and $F_I$ are
ultrametric $K$\nobreakdash-Banach spaces, and $Q^K_{E,E,F_I}$ is analytic at no point. Our argument for this, which we do not
give, is a partition argument with the Erd\H{o}s--Rado theorem in place of Ramsey's theorem. For $k=2$ in
characteristic~$2$ the countable index set $\N$ suffices~\cite{M}.
\end{remark}

\subsection*{Comparison with characteristic \texorpdfstring{$2$}{2}}
For $p=k=2$, Theorem~\ref{thm:main} gives a second counterexample, with a non-ultrametric target, over \emph{every}
field of characteristic~$2$. This includes spherically complete fields such as $\F_2((t))$, where the hypothesis of the
construction of~\cite{M}, that the base field is not spherically complete, fails. Moreover, by
Proposition~\ref{prop:discrete}, which rests on Theorem~A and hence on an unpublished formalization, no counterexample over $\F_2((t))$ can have a complete
ultrametric target such as that of~\cite{M}. Conversely, the target of~\cite{M} is ultrametric. The target of
Theorem~\ref{thm:main} is not, and it has no equivalent ultrametric norm (Proposition~\ref{prop:nonultra}); over
discretely valued bases no complete counterexample target can have one (Proposition~\ref{prop:discrete}). The mechanisms are
different. The characteristic-$2$ argument is a parity identity for determinants (the determinant of an alternating
matrix is the square of its Pfaffian), glued along $\linf$ through the failure of the Hahn--Banach theorem over fields
that are not spherically complete. Here the finite-field theorem is combined with the constant-coefficient map of a
Laurent series field, and Ingleton's theorem is used in the opposite direction, over a spherically complete subfield.
In characteristic~$2$, Theorem~\ref{thm:FF} needs no signs (Remark~\ref{rem:FFhyp}).

\subsection*{Open problems}
All remaining questions concern ultrametric targets.
\begin{enumerate}
\item \emph{Countable index.} Let $K$ be complete, not spherically complete, of characteristic $p\ge3$, and $k\ge p$.
Is there a complete ultrametric $F$ such that $Q^K_{E,E,F}$ is not analytic for $E=\linf(\N,K)$? For $p=2$ the answer
is yes, by~\cite{M}: for $k=2$ this is the theorem of~\cite{M}, and for $k\ge3$ it follows by the wedging argument
of~\cite[Remark~18]{M}, which is formalized in the Lean theorem of \S\ref{sec:lean} under the hypothesis (K2)
of~\cite{M}. The statement announced in Remark~\ref{rem:R} would give such an $F$ only for a much larger index set.
\item \emph{Densely valued spherically complete bases.} Let $K$ be spherically complete, densely valued and of
characteristic $p$, for example the field $\F_p((t^{\mathbb Q}))$ of Hahn series, and let $k\ge p$. Is $Q^K_{E,E',F}$ analytic for all Banach $E,E'$ and every
complete ultrametric $F$? It is when $F$ is spherically complete (Theorem~A), and it is not for some non-ultrametric
Banach $F$ (Theorem~\ref{thm:main}).
\item \emph{Formalization.} Theorem~\ref{thm:main} is not machine-checked. Its most combinatorial part,
Theorem~\ref{thm:FF}, uses only finite combinatorics and Ramsey's theorem, and is a natural first target.
\end{enumerate}

\section{Formalization}\label{sec:lean}
The following statements are formalized in Lean~4 with Mathlib. They compile with no \texttt{sorry}, and the only
axioms they use are \texttt{propext}, \texttt{Quot.\allowbreak{}sound} and \texttt{Classical.\allowbreak{}choice}. Apart from the
pull request~\cite{PR}, the Lean files are not public; the repository of~\cite{M} is private. The Lean names below are
those of this unpublished development (the namespace \texttt{Round24Transfer} is an internal name).
\begin{itemize}\raggedright
\item The $C^n$ statement~\cite{PR}: \texttt{ContinuousAlternatingMap.\allowbreak{}contDiff\_\allowbreak{}compContinuousLinearMapCLM}, for every
$n:\N\cup\{\infty\}$, on the branch of the pull request.
\item The characteristic-$2$ counterexample~\cite{M}, at the point $0$ and $\abs\iota=2$: the theorem
\texttt{AltCounter.\allowbreak{}not\_\allowbreak{}analyticAt\_\allowbreak{}Q\_\allowbreak{}of\_\allowbreak{}NSC} of~\cite{M}, whose hypothesis
\texttt{NSC} is obtained from (K2) by \texttt{AltCounter.\allowbreak{}nsc\_\allowbreak{}of\_\allowbreak{}hasEmptyNestedBalls}. Its extension
to every $\abs\iota\ge2$ and every base point, by the wedging argument of~\cite[Remark~18]{M}:
\texttt{Round24Transfer.\allowbreak{}char2\_\allowbreak{}not\_\allowbreak{}analyticAt\_\allowbreak{}of\_\allowbreak{}hasEmptyNestedBalls}. In both, the base
field is assumed complete, ultrametric and of characteristic~$2$, and to satisfy the hypothesis (K2) of~\cite{M}
(\texttt{AltCounter.\allowbreak{}HasEmptyNestedBalls}): there are closed balls $\cD(c_n,r_n)$ with $\abs{c_{n+1}-c_n}\le r_n$,
$r_0>r_1>\cdots$ and empty intersection. By Remark~\ref{rem:K2}, this is equivalent to the base field not being
spherically complete.
\item The lift criterion, Proposition~\ref{prop:lift}, as the six-fold equivalence \texttt{Round24Transfer.\allowbreak{}tfae\_\allowbreak{}Q}
(bounded lift; continuously polynomial at every point; analytic at every point; analytic at some point; $C^\omega$;
$C^\omega$ at some point), for all normed spaces over any nontrivially normed field (the infinite radius in
Proposition~\ref{prop:lift}(1) is proved on paper).
\item The positive direction of Corollary~\ref{cor:class}:
\texttt{ContinuousAlternatingMap.\allowbreak{}cpolynomialAt\_\allowbreak{}compContinuousLinearMapCLM} and
\texttt{ContinuousAlternatingMap.\allowbreak{}contDiff\_\allowbreak{}compContinuousLinearMapCLM\_\allowbreak{}of\_\allowbreak{}factorial\_\allowbreak{}ne\_\allowbreak{}zero},
under the hypothesis
$((\mathrm{card}\ \iota)!:K)\ne0$, for all normed $E,E',F$.
\item Theorem~A (\S\ref{sec:complements}):
\texttt{ContinuousAlternatingMap.\allowbreak{}contDiff\_\allowbreak{}compContinuousLinearMapCLM\_\allowbreak{}of\_\allowbreak{}sphericallyComplete}.
\item Auxiliary results of this paper: Lemma~\ref{lem:ultra}; the case of Lemma~\ref{lem:discdist} in which $Y$
is a normed field, which gives Lemma~\ref{lem:laurentsub}(4); and a form of Ingleton's theorem
(Theorem~\ref{thm:ingleton}) with a spherically complete ultrametric codomain.
\end{itemize}
Theorem~A and the last two auxiliary facts use a class of spherically complete spaces that is defined in the
unpublished development, with the definition of \S\ref{sec:ascent}; it is not part of Mathlib.
Theorem~\ref{thm:FF}, Proposition~\ref{prop:supp}, Lemma~\ref{lem:Psi}, Theorem~\ref{thm:laurent} and the base change of
\S\ref{sec:ascent} are not formalized, and hence neither are Theorem~\ref{thm:main} and the negative direction of
Corollary~\ref{cor:class}. Proposition~\ref{prop:discrete} is not formalized either, apart from Theorem~A, which it uses.

\begin{thebibliography}{99}
\bibitem{G1} S.~Gou\"ezel, \emph{In characteristic $0$, precomposition by a linear map is analytic on the space of
alternating maps}, pull request \#43338 to Mathlib,\\ {\small\url{https://github.com/leanprover-community/mathlib4/pull/43338}}.
\bibitem{G2} S.~Gou\"ezel, \emph{The bundle of continuous alternating maps is $C^n$}, draft pull request \#43234 to Mathlib,\\
{\small\url{https://github.com/leanprover-community/mathlib4/pull/43234}}.
\bibitem{In} A.~W.~Ingleton, \emph{The Hahn--Banach theorem for non-Archimedean-valued fields}, Proc.\ Cambridge Philos.\
Soc.\ \textbf{48} (1952), 41--45.
\bibitem{La} S.~Lang, \emph{Fundamentals of Differential Geometry}, Graduate Texts in Mathematics 191, Springer, 1999.
\bibitem{docs} The mathlib Community, \emph{Mathlib documentation}, \texttt{docs\#AnalyticAt},\\
{\small\url{https://leanprover-community.github.io/mathlib4_docs/find/?pattern=AnalyticAt#doc}}.
\bibitem{M} J.~McCarthy, \emph{Precomposition on continuous alternating maps is not analytic in characteristic $2$},
note, version of 21 September 2026; Lean formalization in the repository (private)\\
{\small\url{https://github.com/Deicyde/alternating-analytic-counterexample}}.
\bibitem{PR} J.~McCarthy, \emph{Smoothness reflects along a linear isometry with closed range}, draft pull request \#43548
to Mathlib,\\ {\small\url{https://github.com/leanprover-community/mathlib4/pull/43548}}.
\bibitem{Ra} F.~P.~Ramsey, \emph{On a problem of formal logic}, Proc.\ London Math.\ Soc.\ (2) \textbf{30} (1930), 264--286.
\bibitem{Se} J.-P.~Serre, \emph{Endomorphismes compl\`etement continus des espaces de Banach $p$\nobreakdash-adiques}, Publ.\ Math.\
IH\'ES \textbf{12} (1962), 69--85.
\bibitem{vR} A.~C.~M.~van Rooij, \emph{Non-Archimedean Functional Analysis}, Marcel Dekker, 1978.
\end{thebibliography}
\end{document}
